% Chapter 2, Section 2 _Linear Algebra_ Jim Hefferon
%  http://joshua.smcvt.edu/linearalgebra
%  2001-Jun-10
\section{Kebebasan Linear}
Bagian sebelumnya menunjukkan cara memahami suatu ruang vektor
sebagai sebuah rentang,
yaitu sebagai kombinasi linear tanpa pembatasan dari beberapa elemennya.
Sebagai contoh, ruang polinomial linear $\set{a+bx\suchthat a,b\in\Re}$
direntang oleh himpunan $\set{1,x}$.
Bagian sebelumnya juga menunjukkan bahwa suatu ruang dapat memiliki banyak
himpunan yang merentangnya. Dua himpunan lain yang merentang ruang
polinomial linear adalah $\set{1,2x}$ dan $\set{1,x,2x}$.

Pada akhir bagian tersebut, kita menyebut beberapa himpunan perentang sebagai
`minimal', tetapi kita belum mendefinisikan kata itu secara tepat.
% We could mean one of two things.
Yang dapat kita maksud ialah bahwa suatu himpunan perentang bersifat minimal
jika banyak anggotanya paling sedikit di antara semua himpunan dengan rentang
yang sama. Dengan demikian, $\set{1,x,2x}$ tidak minimal karena memiliki tiga
anggota, sedangkan ada himpunan beranggota dua yang merentang ruang yang sama.
Atau, kita dapat memaksudkan bahwa suatu himpunan perentang bersifat minimal
jika tidak ada elemennya yang dapat dibuang tanpa mengubah rentang.
Menurut makna ini, $\set{1,x,2x}$ tidak minimal karena membuang \( 2x \)
sehingga diperoleh \( \set{1,x} \) tidak mengubah rentangnya.

Makna minimalitas yang pertama tampaknya merupakan persyaratan global, sebab
untuk memeriksa apakah suatu himpunan perentang minimal, kita seolah-olah harus
melihat semua himpunan yang merentang dan mencari yang banyak elemennya paling
sedikit. Makna minimalitas yang kedua bersifat lokal karena kita hanya perlu
melihat himpunan itu sendiri dan membandingkan rentangnya ketika berbagai
elemen disertakan atau dibuang. Misalnya, dengan makna kedua, kita dapat
membandingkan rentang $\set{1,x,2x}$ dengan rentang $\set{1,x}$ dan melihat
bahwa $2x$ merupakan suatu ``pengulangan'', dalam arti bahwa membuangnya tidak
memperkecil rentang.

Dalam bagian ini kita akan memakai makna kedua dari `himpunan perentang
minimal' karena kemudahan teknis tersebut. Namun, hasil terpenting dalam buku
ini menyatakan bahwa kedua makna itu berimpit. Kita akan membuktikannya pada
bagian berikutnya.


 






\subsection{Definisi dan Contoh}

Kita telah menjumpai ``pengulangan'' dalam bab pertama. Di sana, Metode Gauss
mengubahnya menjadi persamaan $0=0$.

\begin{example}  \label{ex:StaticsLIAndLD}
Ingat kembali contoh \hyperlink{ex:Statics}{Statika} pada pembukaan Bab Satu.
Kita memperoleh dua keadaan setimbang dengan sepasang benda bermassa tak
diketahui: yang satu pada \( 40 \)~cm dan \( 15 \)~cm, dan yang lain pada
\( -50 \)~cm dan \( 25 \)~cm; lalu kita menghitung nilai kedua massa tersebut.
Seandainya keadaan setimbang kedua justru terjadi pada \( 20 \)~cm dan
\( 7.5 \)~cm, Metode Gauss pada sistem dua persamaan dengan dua variabel yang
dihasilkan tidak akan memberikan satu solusi. Metode itu akan menghasilkan
persamaan $0=0$ bersama sebuah persamaan yang memuat variabel bebas.
Secara intuitif, masalahnya ialah \( \rowvec{20 &7.5} \) merupakan setengah
dari $\rowvec{40 &15}$; yaitu, $\rowvec{20 &7.5}$ berada dalam rentang
himpunan $\set{\rowvec{40 &15}}$ sehingga merupakan data yang berulang.
Kita akan mencoba menyelesaikan masalah dua variabel dengan, pada hakikatnya,
hanya satu informasi.
\end{example}
 
Kita menyebut $\vec{v}$ sebagai suatu ``pengulangan'' dari vektor-vektor dalam
himpunan~$S$ jika $\vec{v}\in\spanof{S}$, sehingga vektor itu bergantung pada,
yakni dapat dinyatakan dalam suku-suku dari, elemen-elemen himpunan tersebut,
$\vec{v}=c_1\vec{s}_1+\cdots+c_n\vec{s}_n$.

\begin{lemma} \label{lm:AddVecSpanNoGrowIffVecInSpan}
%<*lm:AddVecSpanNoGrowIffVecInSpan>
Jika $V$ adalah ruang vektor, $S$ merupakan subhimpunan ruang tersebut, dan
$\vec{v}$ merupakan elemen ruang tersebut, maka
$\spanof{S\union\set{\vec{v}}} = \spanof{S}$
berlaku jika dan hanya jika
$\vec{v}\in\spanof{S}$.
%</lm:AddVecSpanNoGrowIffVecInSpan>
\end{lemma}

\begin{proof}
%<*pf:AddVecSpanNoGrowIffVecInSpan0>
Salah satu arah dari pernyataan jika dan hanya jika ini langsung diperoleh:
jika $\vec{v}\notin\spanof{S}$, kedua himpunan tidak sama karena
$\vec{v}\in\spanof{S\union\set{\vec{v}}}$.

Untuk arah lainnya, andaikan $\vec{v}\in\spanof{S}$ sehingga
$\vec{v}=c_1\vec{s}_1+\cdots+c_n\vec{s}_n$ untuk sejumlah skalar~$c_i$ dan
vektor $\vec{s}_i\in S$. Kita akan memakai ketercakupan timbal balik untuk
menunjukkan bahwa himpunan $\spanof{S\union\set{\vec{v}}}$ dan $\spanof{S}$
sama. Ketercakupan $\spanof{S\union\set{\vec{v}}}\supseteq\spanof{S}$ sudah
jelas.
%</pf:AddVecSpanNoGrowIffVecInSpan0>

%<*pf:AddVecSpanNoGrowIffVecInSpan1>
Untuk menunjukkan ketercakupan dalam arah sebaliknya, ambil $\vec{w}$ sebagai
elemen dari~$\spanof{S\union\set{\vec{v}}}$. Maka $\vec{w}$ merupakan
kombinasi linear dari elemen-elemen $S\union\set{\vec{v}}$, yang dapat kita
tulis sebagai
$\vec{w}=c_{n+1}\vec{s}_{n+1}+\cdots+c_{n+k}\vec{s}_{n+k}+c_{n+k+1}\vec{v}$.
(Mungkin beberapa $\vec{s}_i$ dalam persamaan $\vec{w}$ sama dengan beberapa
yang ada dalam persamaan $\vec{v}$, tetapi hal itu tidak menjadi masalah.)
Substitusikan penjabaran $\vec{v}$.
\begin{equation*}
  \vec{w}=c_{n+1}\vec{s}_{n+1}+\cdots+c_{n+k}\vec{s}_{n+k}
            +c_{n+k+1}\cdot(c_1\vec{s}_1+\cdots+c_n\vec{s}_n)
\end{equation*}
Ruas kanan merupakan kombinasi linear dari kombinasi-kombinasi linear
vektor-vektor dalam~$S$. Jadi, $\vec{w}\in\spanof{S}$.
%</pf:AddVecSpanNoGrowIffVecInSpan1>
\end{proof}

Pembahasan pada pembukaan bagian ini menyangkut pembuangan vektor, bukan
penambahannya.

\begin{corollary}  \label{cor:OmitVecNotShrinkSpanIffVecIsDep}
%<*cor:OmitVecNotShrinkSpanIffVecIsDep>
Untuk $\vec{v}\in S$, membuang vektor itu tidak memperkecil rentang jika dan
hanya jika vektor tersebut bergantung pada vektor-vektor lain dalam himpunan.
Artinya,
$\spanof{S}=\spanof{S-\set{\vec{v}}}$
jika dan hanya jika $\vec{v}\in\spanof{S-\set{\vec{v}}}$.
%</cor:OmitVecNotShrinkSpanIffVecIsDep>
\end{corollary}

Jadi, untuk mengetahui apakah membuang sebuah vektor akan memperkecil rentang,
kita perlu mengetahui apakah vektor itu merupakan kombinasi linear dari
vektor-vektor lain dalam himpunan.

\begin{definition}\label{def:LinInd}
%<*df:LinInd>
Dalam sebarang ruang vektor, suatu himpunan vektor disebut
\definend{bebas linear}\index{bebas linear}%
\index{himpunan!bergantung, bebas}
jika tidak satu pun elemennya merupakan kombinasi linear dari elemen-elemen
lain dalam himpunan itu.\footnote{Lihat juga
\nearbyremark{rem:WhyLIUsesMultiset}.}\spacefactor=1000\ 
Jika tidak, himpunan tersebut disebut
\definend{bergantung linear}\index{bergantung linear}.
%</df:LinInd>
\end{definition}
% A multiset subset of a vector space is
% \definend{linearly independent}\index{linearly independent}%
% \index{sets!dependent, independent}
% if none of its elements is a linear combination of the 
% others.\footnote{More information on multisets is in the appendix. 
% Most of the time we won't need the 
% set-multiset distinction and we will
% follow the standard terminology of referring to a linearly independent or 
% dependent `set'.
% \nearbyremark{rem:WhyLIUsesMultiset} explains why 
% the definition requires a multiset.}\spacefactor=1000\ 
% Otherwise it is \definend{linearly dependent}\index{linearly dependent}.

Dengan demikian, himpunan $\set{\vec{s}_0,\ldots,\vec{s}_n}$ bebas jika tidak
ada kesamaan
$\vec{s}_i=c_0\vec{s}_0+\ldots+c_{i-1}\vec{s}_{i-1}+c_{i+1}\vec{s}_{i+1}+\ldots+c_n\vec{s}_n$.
Penggunaan kata `lain' dalam definisi tersebut berarti bahwa penulisan
$\vec{s}_i$ sebagai kombinasi linear melalui $\vec{s}_i=1\cdot\vec{s}_i$
tidak diperhitungkan.

%<*LinInd>
Perhatikan bahwa, walaupun cara menuliskan satu vektor sebagai kombinasi dari
vektor-vektor lain ini
\begin{equation*}
   \vec{s}_0=\lincombo{c}{\vec{s}}
\end{equation*}
secara visual menonjolkan \( \vec{s}_0 \), secara aljabar tidak ada yang
istimewa tentang vektor itu dalam persamaan tersebut. Untuk setiap
\( \vec{s}_i \) dengan koefisien $c_i$ yang tak nol, kita dapat menulis ulang
persamaan untuk mengisolasi \( \vec{s}_i \).
\begin{equation*}
   \vec{s}_i=(1/c_i)\vec{s}_0+\dots
              +(-c_{i-1}/c_i)\vec{s}_{i-1}+(-c_{i+1}/c_i)\vec{s}_{i+1}
              +\dots+(-c_n/c_i)\vec{s}_n
\end{equation*}
Jika kita tidak ingin mengistimewakan vektor mana pun, kita akan mengatakan
bahwa \( \vec{s}_0,\vec{s}_1,\dots,\vec{s}_n \) berada dalam suatu
\definend{relasi linear}\index{relasi linear}%
\index{relasi!linear}
dan menempatkan semua vektor pada ruas yang sama.
%</LinInd>
Hasil berikut menyatakan kembali definisi kebebasan linear dalam bentuk ini.
Inilah cara yang biasanya kita pakai untuk menghitung apakah suatu himpunan
hingga bergantung atau bebas.

\begin{lemma}   \label{le:LDIffANonTrivLinRel}
%<*lm:LDIffANonTrivLinRel>
Suatu subhimpunan \( S \) dari ruang vektor bebas linear jika dan hanya jika,
di antara elemen-elemennya, satu-satunya relasi linear
$ c_1\vec{s}_1+\dots+c_n\vec{s}_n=\zero$
adalah relasi trivial, \( c_1=0,\dots,\,c_n=0 \)
(dengan $\vec{s}_i\neq\vec{s}_j$ apabila $i\neq j$).
%</lm:LDIffANonTrivLinRel>
\end{lemma}

\begin{proof}
%<*pf:LDIffANonTrivLinRel>
Jika \( S \) bebas linear, tidak ada vektor $\vec{s}_i$ yang merupakan
kombinasi linear dari vektor-vektor lain dalam $S$. Karena itu, tidak ada
relasi linear yang beberapa vektor $\vec{s}$-nya memiliki koefisien tak nol.

Jika \( S \) tidak bebas linear, maka suatu \( \vec{s}_i \) merupakan
kombinasi linear
$\vec{s}_i=c_1\vec{s}_1+\dots+c_{i-1}\vec{s}_{i-1}
    +c_{i+1}\vec{s}_{i+1}+\dots+c_n\vec{s}_n$
dari vektor-vektor lain dalam \( S \). Mengurangkan $\vec{s}_i$ dari kedua
ruas menghasilkan suatu relasi yang melibatkan koefisien tak nol, yaitu
\( -1 \) di depan \( \vec{s}_i \).
%</pf:LDIffANonTrivLinRel>
\end{proof}

\begin{example}  \label{ex:StaticsLI}
Dalam ruang vektor yang terdiri atas vektor baris berlebar dua, himpunan
beranggota dua \( \set{ \rowvec{40 &15},\rowvec{-50 &25}} \) bebas linear.
Untuk memeriksanya, ambil
\begin{equation*}
  c_1\cdot\rowvec{40 &15}+c_2\cdot\rowvec{-50 &25}=\rowvec{0 &0}
\end{equation*}
dan selesaikan sistem yang dihasilkan.
\begin{equation*}
  \begin{linsys}{2}
    40c_1  &-  &50c_2  &=  &0  \\
    15c_1  &+  &25c_2  &=  &0  
   \end{linsys}
  \grstep{-(15/40)\rho_1+\rho_2}
  \begin{linsys}{2}
     40c_1  &- &50c_2       &=  &0  \\
            &  &(175/4)c_2  &=  &0  
   \end{linsys}
\end{equation*}
Baik \( c_1 \) maupun \( c_2 \) sama dengan nol. Jadi, satu-satunya relasi
linear antara kedua vektor baris yang diberikan adalah relasi trivial.

Dalam ruang vektor yang sama, himpunan
\( \set{ \rowvec{40 &15},\rowvec{20 &7.5}} \) bergantung linear karena kita
dapat memenuhi
% \begin{equation*}
$c_1\cdot \rowvec{40 &15}+c_2\cdot\rowvec{20 &7.5}=\rowvec{0 &0}$
% \end{equation*}
dengan \( c_1=1 \) dan \( c_2=-2 \).
\end{example}

\begin{example}
Himpunan \( \set{1+x,1-x} \) bebas linear dalam \( \polyspace_2 \), yaitu
ruang polinomial kuadrat berkoefisien real, karena
\begin{equation*}
   0+0x+0x^2
   =
   c_1(1+x)+c_2(1-x)
   =
   (c_1+c_2)+(c_1-c_2)x+0x^2
\end{equation*}
menghasilkan
\begin{equation*}
  \begin{linsys}{2}
    c_1  &+  &c_2  &=  &0  \\
    c_1  &-  &c_2  &=  &0  
   \end{linsys}
  \grstep{-\rho_1+\rho_2}
  \begin{linsys}{2}
     c_1  &+  &c_2  &=  &0  \\
          &   &2c_2 &=  &0
  \end{linsys}
\end{equation*}
sebab polinomial sama hanya jika koefisien-koefisiennya sama. Jadi,
satu-satunya relasi linear antara kedua anggota $\polyspace_2$ ini adalah
relasi trivial.
\end{example}

\begin{remark}
Lema tersebut menetapkan bahwa $\vec{s}_i\neq\vec{s}_j$ apabila $i\neq j$
karena, tentu saja, jika suatu vektor~$\vec{s}$ muncul dua kali, kita dapat
memperoleh relasi tak trivial $c_1\vec{s}_1+\dots+c_n\vec{s}_n=\zero$ dengan
mengambil koefisien-koefisien terkait sebesar $1$ dan~$-1$. Selain itu, jika
suatu vektor muncul lebih dari sekali dalam sebuah ekspresi, kita selalu dapat
menggabungkan koefisien-koefisiennya.

Perhatikan bahwa lema tersebut juga mengizinkan kebalikan dari kemunculan
lebih dari sekali, yakni beberapa vektor dalam~$S$ tidak muncul sama sekali.
Misalnya, jika~$S$ tak hingga, setiap relasi linear hanya melibatkan sejumlah
hingga vektor sehingga banyak vektor dalam $S$ tidak tercantum. Namun,
perhatikan pula bahwa jika~$S$ hingga, bila memudahkan kita dapat mengambil
kombinasi $c_1\vec{s}_1+\dots+c_n\vec{s}_n$ yang memuat setiap vektor dalam
$S$ tepat satu kali. Jika suatu vektor tidak tercantum, kita dapat
menambahkannya dengan menggunakan koefisien~$0$.
\end{remark}

\begin{example}
Baris-baris matriks ini
\begin{equation*}
  A=
  \begin{mat}[r]
    2  &3  &1  &0  \\
    0  &-1 &0  &-2 \\
    0  &0  &0  &1
  \end{mat}
\end{equation*}
membentuk himpunan bebas linear. Hal ini mudah diperiksa untuk kasus ini,
tetapi ingat pula bahwa Lema~Satu.III.\ref{le:EchFormNoLinCombo} menunjukkan
bahwa baris-baris dari sebarang matriks berbentuk eselon membentuk himpunan
bebas linear.
\end{example}

\begin{example}
Dalam \( \Re^3 \), dengan
\begin{equation*}
   \vec{v}_1=\colvec{3 \\ 4 \\ 5}
   \quad
   \vec{v}_2=\colvec{2 \\ 9 \\ 2}
   \quad
   \vec{v}_3=\colvec{4 \\ 18 \\ 4}
\end{equation*}
himpunan \( S=\set{\vec{v}_1,\vec{v}_2,\vec{v}_3} \) bergantung linear karena
berikut ini merupakan suatu relasi
\begin{equation*}
  0\cdot\vec{v}_1
  +2\cdot\vec{v}_2
  -1\cdot\vec{v}_3
  =\zero
\end{equation*}
dengan skalar-skalar yang tidak semuanya nol (kenyataan bahwa beberapa
skalarnya nol tidak menjadi masalah).
\end{example}

% \begin{remark}  \label{rem:WhyLIIsNonTrivLinRel}
Contoh itu menggambarkan mengapa, walaupun \nearbydefinition{def:LinInd}
menyatakan makna kebebasan dengan lebih jelas,
\nearbylemma{le:LDIffANonTrivLinRel} lebih baik untuk perhitungan. Jika
langsung memakai definisi, seseorang yang mencoba menentukan apakah $S$
bebas linear akan mulai dengan menetapkan
\( \vec{v}_1=c_2\vec{v}_2+c_3\vec{v}_3 \)
dan menyimpulkan bahwa tidak ada $c_2$ dan $c_3$ semacam itu. Namun,
mengetahui bahwa vektor pertama tidak bergantung pada kedua vektor lainnya
belum cukup. Orang tersebut masih harus mencoba
\( \vec{v}_2=c_1\vec{v}_1+c_3\vec{v}_3 \) untuk menemukan kebergantungan
$c_1=0$, \( c_3=1/2 \). \nearbylemma{le:LDIffANonTrivLinRel} memperoleh
kesimpulan yang sama hanya dengan satu perhitungan.
% \end{remark}

\begin{example} \label{ex:EmSetLI}
Subhimpunan kosong\index{himpunan!kosong} dari suatu ruang vektor bebas linear.
Tidak ada relasi linear tak trivial di antara anggotanya karena himpunan itu
tidak memiliki anggota.
\end{example}

\begin{example} \label{ex:SetWithZeroVecLD}
Dalam sebarang ruang vektor, setiap subhimpunan yang memuat vektor nol
bergantung linear. Salah satu contohnya, dalam ruang $\polyspace_2$ yang
terdiri atas polinomial kuadrat, ialah subhimpunan $\set{1+x,x+x^2,0}$.
Himpunan itu bergantung linear karena
$0\cdot\vec{v}_1+0\cdot\vec{v}_2+1\cdot\zero=\zero$ merupakan relasi tak
trivial, sebab koefisien-koefisiennya tidak semuanya nol.
\end{example}

Ada satu hal halus yang akan kita jumpai beberapa kali dan berkaitan dengan
contoh sebelumnya. Hal itu menyangkut jumlah trivial, yaitu jumlah himpunan
kosong. Salah satu cara untuk melihat bagaimana jumlah trivial harus
didefinisikan ialah mempertimbangkan urutan
$\vec{v}_1+\vec{v}_2+\vec{v}_3$, diikuti oleh
$\vec{v}_1+\vec{v}_2$, lalu
$\vec{v}_1$.
Selisih antara jumlah tiga vektor dan jumlah dua vektor adalah $\vec{v}_3$.
Kemudian, selisih antara jumlah dua vektor dan jumlah satu vektor adalah
$\vec{v}_2$. Ketika berikutnya beralih ke jumlah trivial, yaitu jumlah tanpa
vektor, kita dapat mengharapkan pengurangan~$\vec{v}_1$. Karena itu, kita
mendefinisikan jumlah tanpa vektor sebagai vektor nol.

Kaitannya dengan contoh sebelumnya ialah bahwa jika vektor nol berada dalam
suatu himpunan, himpunan tersebut memiliki elemen yang merupakan kombinasi dari
suatu subhimpunan vektor-vektor lainnya; khususnya, vektor nol merupakan
kombinasi dari subhimpunan kosong. Bahkan himpunan $S=\set{\zero}$ bergantung
linear, karena $\zero$ adalah jumlah himpunan kosong dan himpunan kosong
merupakan subhimpunan dari~$S$.

\begin{remark} \label{rem:WhyLIUsesMultiset} %\cite{Velleman} 
Definisi kebebasan linear, \nearbydefinition{def:LinInd}, merujuk pada suatu
`himpunan' vektor. Himpunan merupakan jenis koleksi yang paling dikenal dan,
dalam praktik, semua orang memakai kata `himpunan' dalam konteks ini. Namun,
demi kelengkapan, kita akan mencatat bahwa himpunan bukanlah jenis koleksi yang
sepenuhnya tepat untuk tujuan ini.

Ingat bahwa suatu himpunan adalah koleksi dengan dua sifat:
(i)~urutan tidak berpengaruh, sehingga himpunan $\set{1,2}$ sama dengan
himpunan $\set{2,1}$, dan (ii)~duplikat melebur, sehingga himpunan
$\set{1,1,2}$ sama dengan himpunan $\set{1,2}$.

Sekarang perhatikan reduksi matriks ini.
\begin{equation*}
  \begin{mat}
    1 &1 &1 \\
    2 &2 &2 \\
    1 &2 &3
  \end{mat}
  \grstep{(1/2)\rho_2}
  \begin{mat}
    1 &1 &1 \\
    1 &1 &1 \\
    1 &2 &3
  \end{mat}
\end{equation*}
Di sebelah kiri, himpunan baris matriks
$\set{\rowvec{1 &1 &1}, \rowvec{2 &2 &2}, \rowvec{1 &2 &3}}$ bergantung
linear. Di sebelah kanan, himpunan barisnya adalah
$\set{\rowvec{1 &1 &1}, \rowvec{1 &1 &1}, \rowvec{1 &2 &3}}$.
Karena duplikat melebur, himpunan itu sama dengan
$\set{\rowvec{1 &1 &1}, \rowvec{1 &2 &3}}$, yang bebas linear. Ini menjadi
masalah karena Metode Gauss seharusnya mempertahankan kebergantungan linear.

Artinya, secara tegas, kita memerlukan jenis koleksi yang duplikatnya tidak
melebur. Koleksi yang urutannya tidak berpengaruh dan duplikatnya tidak melebur
disebut \definend{multihimpunan}.\index{himpunan!multihimpunan}%
\index{multihimpunan}\index{kebebasan linear!multihimpunan}

Namun, meskipun bersikeras untuk sepenuhnya tepat memiliki keuntungan,
menyimpang dari istilah baku `himpunan' juga menimbulkan kesulitannya sendiri.
Karena itu, kita akan terus memakai kata tersebut. Kelak, kadang-kadang kita
perlu membentuk kombinasi tanpa membiarkan duplikat melebur, dan kita akan
melakukannya tanpa komentar lebih lanjut.
\end{remark}

\begin{corollary} \label{cor:LiSetMinSpanSet}
%<*cor:LiSetMinSpanSet>
Suatu himpunan $S$ bebas linear jika dan hanya jika, untuk setiap
$\vec{v}\in S$, pembuangannya memperkecil rentang
$\spanof{S-\set{\vec{v}}}\subsetneq\spanof{S}$.
%</cor:LiSetMinSpanSet>
\end{corollary}

\begin{proof}
%<*pf:LiSetMinSpanSet>
Ini mengikuti \nearbycorollary{cor:OmitVecNotShrinkSpanIffVecIsDep}. Jika $S$
bebas linear, tidak ada vektornya yang bergantung pada elemen-elemen lain,
sehingga pembuangan sebarang vektor akan memperkecil rentang. Jika $S$ tidak
bebas linear, himpunan itu memuat sebuah vektor yang bergantung pada
elemen-elemen lain dalam himpunan, dan pembuangan vektor tersebut tidak akan
memperkecil rentang.
%</pf:LiSetMinSpanSet>
\end{proof}

Jadi, suatu himpunan perentang minimal jika dan hanya jika bebas linear.

Hasil sebelumnya membahas pembuangan elemen dari himpunan bebas linear. Hasil
berikutnya membahas penambahan elemen.

\begin{lemma}  \label{lm:AddVecLiSetIsLiIffVecNotInSpan}
%<*lm:AddVecLiSetIsLiIffVecNotInSpan>
Andaikan $S$ bebas linear dan $\vec{v}\notin S$. Maka himpunan
$S\union\set{\vec{v}}$ bebas linear jika dan hanya jika
$\vec{v}\notin\spanof{S}$.   
%</lm:AddVecLiSetIsLiIffVecNotInSpan>
\end{lemma}

\begin{proof}
%<*pf:AddVecLiSetIsLiIffVecNotInSpan0>
% Assume that $S$ is linearly independent and that $\vec{v}\notin S$.
Kita akan menunjukkan bahwa $S\union\set{\vec{v}}$ tidak bebas linear jika dan
hanya jika $\vec{v}\in\spanof{S}$.

Pertama, andaikan $\vec{v}\in\spanof{S}$. Nyatakan $\vec{v}$ sebagai kombinasi
$\vec{v}=c_1\vec{s}_1+\cdots+c_n\vec{s}_n$. Tulis ulang sebagai
$\zero=c_1\vec{s}_1+\cdots+c_n\vec{s}_n-1\cdot\vec{v}$. Karena
$\vec{v}\notin S$, vektor itu tidak sama dengan satu pun $\vec{s}_i$,
sehingga ini merupakan kebergantungan linear tak trivial di antara elemen
$S\union\set{\vec{v}}$. Jadi, himpunan itu tidak bebas linear.
%</pf:AddVecLiSetIsLiIffVecNotInSpan0>

%<*pf:AddVecLiSetIsLiIffVecNotInSpan1>
Sekarang andaikan $S\union\set{\vec{v}}$ tidak bebas linear dan perhatikan
suatu kebergantungan tak trivial di antara anggotanya,
$\zero=c_1\vec{s}_1+\cdots+c_n\vec{s}_n+c_{n+1}\cdot\vec{v}$.
Jika $c_{n+1}=0$, ini merupakan kebergantungan di antara elemen-elemen~$S$.
Namun, kita mengandaikan bahwa~$S$ bebas, sehingga $c_{n+1}\neq 0$. Tulis ulang
persamaan tersebut sebagai
$\vec{v}=(c_1/c_{n+1})\vec{s}_1+\cdots+(c_n/c_{n+1})\vec{s}_n$ untuk
memperoleh $\vec{v}\in\spanof{S}$.
%</pf:AddVecLiSetIsLiIffVecNotInSpan1>
\end{proof}



\begin{example} \label{ex:LinindSetsAndSuper}
Subhimpunan dari $\Re^3$ ini bebas linear.
\begin{equation*}
  S
   =\set{\colvec{1 \\ 0 \\ 0}}
\end{equation*}
Rentang $S$ adalah sumbu-$x$. Berikut dua himpunan super, yang satu bergantung
linear dan yang lain bebas.
\begin{center}
     bergantung:
     \( \set{
         \colvec{1 \\ 0 \\ 0},
         \colvec{-3 \\ 0 \\ 0}} \)      
     \qquad
     bebas:
     \( \set{
         \colvec{1 \\ 0 \\ 0},
         \colvec{0 \\ 1 \\ 0}} \)      
\end{center}
Kita memperoleh himpunan super yang bergantung dengan menambahkan vektor dari
sumbu-$x$, sehingga rentangnya tidak membesar. Kita memperoleh himpunan super
yang bebas dengan menambahkan vektor yang tidak berada dalam $\spanof{S}$,
karena vektor itu memiliki komponen~$y$ tak nol, sehingga rentangnya membesar.

Untuk himpunan bebas
\begin{equation*}
  S
   =\set{\colvec{1 \\ 0 \\ 0},
          \colvec{0 \\ 1 \\ 0}
                 }
\end{equation*}
rentang \( \spanof{S} \) adalah bidang-\( xy \). Berikut dua himpunan super.
\begin{center}
     bergantung:
     \( \set{
         \colvec{1 \\ 0 \\ 0},
         \colvec{0 \\ 1 \\ 0},
         \colvec{3 \\ -2 \\ 0} } \)       
     \qquad
     bebas:
     \( \set{
         \colvec{1 \\ 0 \\ 0},
         \colvec{0 \\ 1 \\ 0},
         \colvec{0 \\ 0 \\ 1} } \)    
\end{center}
Seperti di atas, anggota tambahan pada himpunan super yang bergantung berasal
dari $\spanof{S}$, yaitu bidang-$xy$, sedangkan anggota tambahan pada himpunan
super yang bebas berasal dari luar rentang tersebut.

Terakhir, perhatikan himpunan bebas ini
\begin{equation*}
   S =\set{\colvec{1 \\ 0 \\ 0},
          \colvec{0 \\ 1 \\ 0},
          \colvec{0 \\ 0 \\ 1}        } 
\end{equation*}
dengan \( \spanof{S}=\Re^3 \). Kita dapat memperoleh himpunan super yang
bergantung linear.
\begin{center}
     bergantung:
     \( \set{
         \colvec{1 \\ 0 \\ 0},
         \colvec{0 \\ 1 \\ 0},
         \colvec{0 \\ 0 \\ 1},
         \colvec{2 \\ -1 \\ 3} } \)     
\end{center}
Namun, tidak ada himpunan super bebas linear dari $S$. Salah satu cara untuk
melihatnya ialah mencatat bahwa, untuk sebarang vektor yang hendak kita
tambahkan ke $S$, persamaan
\begin{equation*}
  \colvec{x \\ y \\ z}
  =c_1\colvec{1 \\ 0 \\ 0}
   +c_2\colvec{0 \\ 1 \\ 0}
   +c_3\colvec{0 \\ 0 \\ 1}
\end{equation*}
memiliki solusi $c_1=x$, $c_2=y$, dan $c_3=z$. Cara lain untuk melihatnya ialah
bahwa kita tidak dapat menambahkan vektor apa pun dari luar rentang
$\spanof{S}$ karena rentang itu adalah $\Re^3$.
\end{example}

\begin{corollary} \label{th:AlwaysAnLDSubset}
%<*th:AlwaysAnLDSubset>
Dalam suatu ruang vektor, setiap himpunan hingga memiliki subhimpunan bebas
linear dengan rentang yang sama.
%</th:AlwaysAnLDSubset>
\end{corollary}

\begin{proof}
%<*pf:AlwaysAnLDSubset0>
Jika \( S=\set{ \vec{s}_1,\dots,\vec{s}_n} \) bebas linear, maka $S$ sendiri
memenuhi pernyataan tersebut. Jadi, andaikan $S$ bergantung linear.
%</pf:AlwaysAnLDSubset0>

%<*pf:AlwaysAnLDSubset1>
Berdasarkan definisi kebergantungan, $S$ memuat sebuah vektor $\vec{v}_1$ yang
merupakan kombinasi linear dari vektor-vektor lainnya. Definisikan himpunan
\( S_1=S-\set{\vec{v}_1} \). Menurut
\nearbycorollary{cor:OmitVecNotShrinkSpanIffVecIsDep}, % {lm:ShrinkSpanByRemovingNonRepeat} 
rentangnya tidak mengecil: \( \spanof{S_1}=\spanof{S} \).
%</pf:AlwaysAnLDSubset1>

%<*pf:AlwaysAnLDSubset2>
Jika \( S_1 \) bebas linear, pembuktian selesai. Jika tidak, ulangi prosesnya:
ambil sebuah vektor $\vec{v}_2$ yang merupakan kombinasi linear dari anggota
lain $S_1$, lalu buang vektor itu untuk memperoleh
\( S_2=S_1-\set{\vec{v}_2} \) sedemikian sehingga
\( \spanof{S_2}=\spanof{S_1} \). Ulangi sampai muncul suatu himpunan bebas
linear $S_j$; himpunan itu pada akhirnya pasti muncul karena \( S \) hingga
dan himpunan kosong bebas linear.
%</pf:AlwaysAnLDSubset2>
(Secara formal, argumen ini memakai induksi pada banyak elemen dalam $S$.
\nearbyexercise{exer:FillIndDetProofSetHasLISub} meminta rinciannya.)
\end{proof}

Jadi, jika kita memiliki himpunan yang bergantung linear, kita dapat
merampingkannya tanpa mengubah rentang dengan membuang vektor-vektor yang kita
sebut ``pengulangan''.

\begin{example}  \label{ex:ShrinkSetSameSpan}
Himpunan ini merentang \( \Re^3 \) (pemeriksaannya rutin), tetapi tidak bebas
linear.
\begin{equation*}
  S=\set{\colvec[r]{1 \\ 0 \\ 0},
     \colvec[r]{0 \\ 2 \\ 0},
     \colvec[r]{1 \\ 2 \\ 0},
     \colvec[r]{0 \\ -1 \\ 1},
     \colvec[r]{3 \\ 3 \\ 0}  }
\end{equation*}
Kita akan menghitung vektor mana yang harus dibuang untuk memperoleh
subhimpunan yang bebas tetapi memiliki rentang yang sama. Relasi linear ini
\begin{equation*}
  c_1\colvec[r]{1 \\ 0 \\ 0}
  +c_2\colvec[r]{0 \\ 2 \\ 0}
  +c_3\colvec[r]{1 \\ 2 \\ 0}
  +c_4\colvec[r]{0 \\ -1 \\ 1}
  +c_5\colvec[r]{3 \\ 3 \\ 0}
  =\colvec[r]{0 \\ 0 \\ 0}
  \qquad\tag{$*$}
\end{equation*}
menghasilkan suatu sistem
\begin{equation*}
  \begin{linsys}{5}
     c_1  &   &      &+  &c_3   &+  &    &+  &3c_5 &= &0  \\
          &   &2c_2  &+  &2c_3  &-  &c_4 &+  &3c_5 &= &0  \\
          &   &      &   &     &   &c_4  &   &     &= &0  
\end{linsys}
\end{equation*}
yang himpunan solusinya memiliki parametrisasi berikut.
\begin{equation*}
  \set{\colvec{c_1 \\ c_2 \\ c_3 \\ c_4 \\ c_5}=
     c_3\colvec[r]{-1 \\ -1 \\ 1 \\ 0 \\ 0}
     +c_5\colvec[r]{-3 \\ -3/2 \\ 0 \\ 0 \\ 1}
     \suchthat c_3,c_5\in\Re }
\end{equation*}
Tetapkan $c_5=1$ dan~$c_3=0$ untuk memperoleh satu contoh dari ($*$).
\begin{equation*}
  -3\cdot\colvec[r]{1 \\ 0 \\ 0}
  -\frac{3}{2}\cdot\colvec[r]{0 \\ 2 \\ 0}
  +0\cdot\colvec[r]{1 \\ 2 \\ 0}
  +0\cdot\colvec[r]{0 \\ -1 \\ 1}
  +1\cdot\colvec[r]{3 \\ 3 \\ 0}
  =\colvec[r]{0 \\ 0 \\ 0}
\end{equation*}
Ini menunjukkan bahwa vektor dalam $S$ yang kita kaitkan dengan~$c_5$ berada
dalam rentang himpunan yang terdiri atas vektor milik $c_1$ dan vektor milik
$c_2$. Kita dapat membuang vektor kelima $S$ tanpa memperkecil rentang.

Demikian pula, tetapkan $c_3=1$ dan $c_5=0$ untuk memperoleh satu contoh
dari~($*$) yang menunjukkan bahwa kita dapat membuang vektor ketiga $S$ tanpa
memperkecil rentang. Jadi, himpunan berikut memiliki rentang yang sama dengan
$S$.
\begin{equation*}
  % \hat{S}=
     \set{\colvec[r]{1 \\ 0 \\ 0},
     \colvec[r]{0 \\ 2 \\ 0},
     \colvec[r]{0 \\ -1 \\ 1}  }
\end{equation*}
Pemeriksaan bahwa himpunan ini bebas linear bersifat rutin.
\end{example}

\begin{corollary} \label{cor:LDMeansLC}
%<*co:LDMeansLC>
Suatu subhimpunan \( S=\set{\vec{s}_1,\dots,\vec{s}_n} \) dari ruang vektor
bergantung linear jika dan hanya jika suatu \( \vec{s}_i \) merupakan
kombinasi linear dari vektor-vektor \( \vec{s}_1 \), \ldots,
\( \vec{s}_{i-1} \) yang tercantum sebelumnya.
%</co:LDMeansLC>
\end{corollary}

\begin{proof}
%<*pf:LDMeansLC>
Perhatikan \( S_0=\set{} \), \( S_1=\set{\vec{s}_1} \),
\( S_2=\set{\vec{s}_1,\vec{s}_2 } \), dan seterusnya. Suatu indeks
\( i\geq 1 \) merupakan indeks pertama yang membuat
\( S_{i-1}\union\set{\vec{s}_i } \) bergantung linear, dan pada titik itu
\( \vec{s}_i\in\spanof{ S_{i-1} } \).
%</pf:LDMeansLC>
\end{proof}

Pembuktian \nearbycorollary{th:AlwaysAnLDSubset} menjelaskan cara menghasilkan
himpunan bebas linear dengan merampingkan, yakni dengan mengambil
subhimpunan. Pembuktian \nearbycorollary{cor:LDMeansLC} menjelaskan cara
menemukan himpunan yang bergantung linear dengan mengambil himpunan super.
Kita mengakhiri subbagian ini dengan meninjau bagaimana kebebasan dan
kebergantungan linear berinteraksi dengan relasi subhimpunan antarhimpunan.

\begin{lemma}  \label{le:SubsetPreserveLI}
%<*lm:SubsetPreserveLI>
Setiap subhimpunan dari himpunan bebas linear juga bebas linear. Setiap
himpunan super dari himpunan yang bergantung linear juga bergantung linear.
%</lm:SubsetPreserveLI>
\end{lemma}

\begin{proof}
%<*pf:SubsetPreserveLI>
Keduanya jelas.
%</pf:SubsetPreserveLI>
\end{proof}

%<*SubsetPreserveLI>
Dengan kata lain, pengambilan subhimpunan mempertahankan kebebasan dan
pengambilan himpunan super mempertahankan kebergantungan.
%</SubsetPreserveLI>

Itulah dua dari empat kemungkinan kasus. Kasus ketiga, apakah pengambilan
subhimpunan mempertahankan kebergantungan linear, dibahas dalam
\nearbyexample{ex:ShrinkSetSameSpan}, yang memberikan suatu himpunan
bergantung linear $S$ dengan satu subhimpunan yang bergantung linear dan
subhimpunan lain yang bebas. Kasus keempat, apakah pengambilan himpunan super
mempertahankan kebebasan linear, dibahas dalam
\nearbyexample{ex:LinindSetsAndSuper}, yang memberikan kasus ketika suatu
himpunan bebas linear memiliki himpunan super yang bebas maupun yang
bergantung. Tabel berikut merangkumnya.
\smallskip
%<*IndependenceAndSubsetTable>
\begin{center} \renewcommand{\arraystretch}{1.1}
  \begin{tabular}[b]{r|cc} 
                        \multicolumn{1}{c}{}
                        &\multicolumn{1}{c}{\( \hat{S}\subset S \)}
                        &\multicolumn{1}{c}{\( \hat{S}\supset S \)}      \\
     \cline{2-3}\rule{0em}{12pt} % make box a bit taller
          \textit{$S$ bebas} 
              &\( \hat{S} \) harus bebas   &\( \hat{S} \) dapat bebas atau bergantung\\
     % \cline{2-3}\noalign{\smallskip}
          \textit{$S$ bergantung} 
              &\( \hat{S} \) dapat bebas atau bergantung &\( \hat{S} \) harus bergantung    \\
     % \cline{2-3}
   \end{tabular}
\end{center}
%</IndependenceAndSubsetTable>
\smallskip

\nearbyexample{ex:LinindSetsAndSuper} juga memberi tahu kita hal lain tentang
interaksi antara kebebasan linear dan himpunan super. Contoh itu menyebut suatu
himpunan bebas linear yang maksimal, dalam arti tidak memiliki himpunan super
yang bebas linear. Berdasarkan \nearbylemma{lm:AddVecLiSetIsLiIffVecNotInSpan},
suatu himpunan bebas linear maksimal jika dan hanya jika merentang seluruh
ruang, sebab tepat pada saat itulah semua vektor dalam ruang sudah berada
dalam rentang. Hal ini melengkapi \nearbycorollary{cor:LiSetMinSpanSet} dengan
indah: suatu himpunan perentang minimal jika dan hanya jika bebas linear.

% In summary,
% we have introduced the definition of linear independence to 
% formalize the idea of the minimality of a spanning set.
% We have developed some properties of this idea.
% The most important is \nearbylemma{lm:AddVecLiSetIsLiIffVecNotInSpan}, which 
% tells us that a linearly independent set is maximal when it spans the space.


\begin{exercises}
  \recommended \item
    Tentukan apakah setiap subhimpunan \( \Re^3 \) berikut bergantung linear
    atau bebas linear.
    \begin{exparts*}
      \partsitem \( \set{\colvec{1 \\ -3 \\ 5},
                    \colvec{2 \\ 2 \\ 4},
                    \colvec{4 \\ -4 \\ 14} }  \)
      \partsitem \( \set{\colvec{1 \\ 7 \\ 7},
                    \colvec{2 \\ 7 \\ 7},
                    \colvec{3 \\ 7 \\ 7} }  \)
      \partsitem \( \set{\colvec{0 \\ 0 \\ -1},
                    \colvec{1 \\ 0 \\ 4} }  \)
      \partsitem \( \set{\colvec{9 \\ 9 \\ 0},
                    \colvec{2 \\ 0 \\ 1},
                    \colvec{3 \\ 5 \\ -4},
                    \colvec{12 \\ 12 \\ -1} }  \)
    \end{exparts*}
    \begin{answer}
      Untuk masing-masing subhimpunan ini, jika subhimpunan tersebut bebas,
      Anda harus membuktikannya; jika bergantung, Anda harus memberikan contoh
      suatu kebergantungan.
      \begin{exparts}
        \partsitem Himpunan ini bergantung.
          Dengan memperhatikan
          \begin{equation*}
             c_1\colvec{1 \\ -3 \\ 5}
             +c_2\colvec{2 \\ 2 \\ 4}
             +c_3\colvec{4 \\ -4 \\ 14}
             =\colvec{0 \\ 0 \\ 0}
          \end{equation*}
          diperoleh sistem linear berikut.
          \begin{equation*}
            \begin{linsys}{3}
              c_1  &+  &2c_2  &+  &4c_3  &=  &0  \\
              -3c_1&+  &2c_2  &-  &4c_3  &=  &0  \\
              5c_1 &+  &4c_2  &+  &14c_3 &=  &0  
            \end{linsys}
          \end{equation*}
          Metode Gauss
          \begin{equation*}
            \begin{amat}[r]{3}
              1  &2  &4  &0  \\
              -3 &2  &-4 &0  \\
              5  &4  &14 &0
            \end{amat}
            \grstep[-5\rho_1+\rho_3]{3\rho_1+\rho_2}
            \repeatedgrstep{(3/4)\rho_2+\rho_3}            
            \begin{amat}[r]{3}
              1  &2  &4  &0  \\
              0  &8  &8  &0  \\
              0  &0  &0  &0
            \end{amat}
          \end{equation*}
          menghasilkan sebuah variabel bebas, sehingga ada tak hingga banyak
          solusi. Untuk contoh suatu kebergantungan tertentu, kita dapat
          menetapkan $c_3$, misalnya, sebesar $1$. Kemudian kita memperoleh
          \( c_2=-1 \) dan \( c_1=-2 \).
        \partsitem Himpunan ini bergantung.
          Sistem linear yang muncul di sini
          \begin{equation*}
            \begin{amat}[r]{3}
              1  &2  &3  &0  \\
              7  &7  &7  &0  \\
              7  &7  &7  &0
            \end{amat}
            \grstep[-7\rho_1+\rho_3]{-7\rho_1+\rho_2}
            \repeatedgrstep{-\rho_2+\rho_3}
            \begin{amat}[r]{3}
              1  &2  &3   &0  \\
              0  &-7 &-14 &0  \\
              0  &0  &0   &0
            \end{amat}
          \end{equation*}
          memiliki tak hingga banyak solusi. Kita dapat memperoleh solusi
          tertentu dengan mengambil $c_3$, misalnya, sebesar $1$, lalu
          melakukan substitusi balik untuk mendapatkan $c_2$ dan $c_1$.
        \partsitem Himpunan ini bebas linear.
          Sistem
          \begin{equation*}
            \begin{amat}[r]{2}
              0  &1  &0  \\
              0  &0  &0  \\
              -1 &4  &0
            \end{amat}
            \grstep{\rho_1\leftrightarrow\rho_2}
            \repeatedgrstep{\rho_3\leftrightarrow\rho_1}
            \begin{amat}{2}
              -1 &4  &0  \\
              0  &1  &0  \\
              0  &0  &0  
            \end{amat}
          \end{equation*}
          hanya memiliki solusi $c_1=0$ dan $c_2=0$. (Kita juga dapat
          memperoleh jawaban melalui pengamatan\Dash vektor kedua jelas bukan
          kelipatan vektor pertama, dan sebaliknya.)
        \partsitem Himpunan ini bergantung linear.
          Sistem linear
          \begin{equation*}
            \begin{amat}[r]{4}
              9  &2  &3  &12  &0  \\
              9  &0  &5  &12  &0  \\
              0  &1  &-4 &-1  &0
            \end{amat}
          \end{equation*}
          memiliki lebih banyak variabel tak diketahui daripada persamaan,
          sehingga Metode Gauss harus berakhir dengan setidaknya satu variabel
          bebas (tidak mungkin ada persamaan yang kontradiktif karena sistem
          ini homogen, sehingga setidaknya memiliki solusi yang semuanya nol).
          Untuk memperlihatkan suatu kombinasi, kita dapat melakukan reduksi
          \begin{equation*}
            \grstep{-\rho_1+\rho_2}
            \grstep{(1/2)\rho_2+\rho_3}
            \begin{amat}[r]{4}
              9  &2  &3  &12  &0  \\
              0  &-2 &2  &0   &0  \\
              0  &0  &-3 &-1  &0
            \end{amat}
          \end{equation*}
          dan mengambil, misalnya, $c_4=1$. Maka kita memperoleh
          $c_3=-1/3$, $c_2=-1/3$, dan $c_1=-31/27$.
      \end{exparts}  
    \end{answer}
  \recommended \item 
    Manakah dari subhimpunan \( \polyspace_3 \) berikut yang bergantung linear
    dan manakah yang bebas?
    \begin{exparts}
      \partsitem \( \set{3-x+9x^2,5-6x+3x^2,1+1x-5x^2} \)
      \partsitem \( \set{-x^2,1+4x^2} \)
      \partsitem \( \set{2+x+7x^2,3-x+2x^2,4-3x^2} \)
      \partsitem \( \set{8+3x+3x^2,x+2x^2,2+2x+2x^2,8-2x+5x^2} \)
    \end{exparts}
    \begin{answer}
      Dalam kasus kebebasan, Anda harus membuktikan bahwa himpunannya bebas.
      Jika tidak, Anda harus memperlihatkan suatu kebergantungan. (Di sini kami
      memberikan satu kebergantungan tertentu, tetapi yang lain juga mungkin.)
      \begin{exparts}
        \partsitem Himpunan ini bebas.
          Menyusun relasi
          \( c_1(3-x+9x^2)+c_2(5-6x+3x^2)+c_3(1+1x-5x^2)=0+0x+0x^2 \)
          menghasilkan suatu sistem linear
          \begin{equation*}
            \begin{amat}[r]{3}
              3  &5  &1  &0  \\
              -1 &-6 &1  &0  \\
              9  &3  &-5 &0  
            \end{amat}
            \grstep[-3\rho_1+\rho_3]{(1/3)\rho_1+\rho_2}
            \repeatedgrstep{3\rho_2}
            \repeatedgrstep{-(12/13)\rho_2+\rho_3}
            \begin{amat}[r]{3}
              3  &5   &1        &0  \\
              0  &-13 &4        &0  \\
              0  &0   &-128/13  &0  
            \end{amat}
          \end{equation*}
          yang hanya memiliki satu solusi: \( c_1=0 \), \( c_2=0 \), dan
          \( c_3=0 \).
        \partsitem Himpunan ini bebas.
           Kita dapat melihatnya langsung melalui pengamatan berdasarkan
           definisi kebebasan linear. Jelas bahwa yang satu bukan kelipatan
           dari yang lain.
        \partsitem Himpunan ini bebas linear.
           Sistem linearnya direduksi sebagai berikut
           \begin{equation*}
             \begin{amat}[r]{3}
               2  &3  &4  &0  \\
               1  &-1 &0  &0  \\
               7  &2  &-3 &0  
             \end{amat}
             \grstep[-(7/2)\rho_1+\rho_3]{-(1/2)\rho_1+\rho_2}
             \repeatedgrstep{-(17/5)\rho_2+\rho_3}
             \begin{amat}[r]{3}
               2  &3    &4      &0  \\
               0  &-5/2 &-2     &0  \\
               0  &0    &-51/5  &0  
             \end{amat}
           \end{equation*}
           dan menunjukkan bahwa satu-satunya solusi ialah $c_1=0$, $c_2=0$,
           dan $c_3=0$.
        \partsitem Himpunan ini bergantung linear.
           Sistem linear
           \begin{equation*}
             \begin{amat}[r]{4}
               8  &0  &2  &8  &0  \\ 
               3  &1  &2  &-2 &0  \\
               3  &2  &2  &5  &0
             \end{amat}
           \end{equation*}
           setelah direduksi harus berakhir dengan setidaknya satu variabel
           bebas (ada lebih banyak variabel daripada persamaan, dan persamaan
           yang kontradiktif tidak mungkin muncul karena sistemnya homogen).
           Kita dapat mengambil variabel bebas sebagai parameter untuk
           mendeskripsikan himpunan solusi. Kemudian kita dapat menetapkan
           parameter itu pada nilai tak nol untuk memperoleh relasi linear tak
           trivial.
      \end{exparts}  
     \end{answer}
  \item Tentukan apakah setiap himpunan bebas linear dalam ruang alaminya.
    \begin{exparts*}
      \partsitem $\set{\colvec{1 \\ 2 \\ 0}, 
              \colvec{-1 \\ 1 \\ 0}}$
      \partsitem $\set{\rowvec{1 &3 &1}, \rowvec{-1 &4 &3}, \rowvec{-1 &11 &7}}$
      \partsitem $\set{\begin{mat}
                5 &4 \\
                1 &2 
              \end{mat},
              \begin{mat}
                0 &0 \\
                0 &0
              \end{mat},
              \begin{mat}
                1 &0 \\
               -1 &4
              \end{mat}
           }$
    \end{exparts*}
    \begin{answer}
      \begin{exparts}
        \partsitem Ruang vektor alaminya adalah $\Re^3$.
          Susun persamaan
          \begin{equation*}
            \colvec{0 \\ 0 \\ 0}=c_1\colvec{1 \\ 2 \\ 0}
                                 +c_2\colvec{-1 \\ 1 \\ 0}
          \end{equation*}
          dan perhatikan sistem homogen yang dihasilkan.
          \begin{equation*}
             \begin{amat}{2}
               1  &-1 &0  \\
               2  &1  &0  \\
               0  &0  &0
             \end{amat}
             \grstep{-2\rho_1+\rho_2}
             \begin{amat}{2}
              1  &-1 &0  \\
              0  &3  &0  \\
              0  &0  &0
            \end{amat} 
          \end{equation*}
          Sistem ini memiliki solusi tunggal, yaitu $c_1=0$, $c_2=0$. Jadi,
          himpunannya bebas linear.
        \partsitem Ruang vektor alaminya adalah himpunan vektor baris
          berlebar tiga. Persamaan
          \begin{equation*}
            \rowvec{0 &0 &0}=c_1\rowvec{1 &3 &1}
                            +c_2\rowvec{-1 &4 &3}
                            +c_3\rowvec{-1 &11 &7}
          \end{equation*}
          menghasilkan suatu sistem linear
          \begin{align*}
            \begin{amat}{3}
              1  &-1  &-1  &0  \\
              3  &4   &11  &0  \\
              1  &3   &7   &0
           \end{amat}
           &\grstep[-\rho_1+\rho_3]{-3\rho_1+\rho_2}
           \begin{amat}{3}
             1  &-1  &-1  &0  \\
             0  &7   &14  &0  \\
             0  &4   &8   &0
          \end{amat}                                   \\
          &\grstep{(4/7)\rho_2+\rho_3}
          \begin{amat}{3}
            1  &-1  &-1  &0  \\
            0  &7   &14  &0  \\
            0  &0   &0   &0
         \end{amat}
       \end{align*}
       dengan tak hingga banyak solusi, yaitu lebih dari sekadar solusi
       trivial.
       \begin{equation*}
         \set{\colvec{c_1 \\ c_2 \\ c_3}
              =\colvec{-1 \\ -2 \\ 1}c_3
              \suchthat c_3\in\Re}
       \end{equation*}
       Jadi, himpunan itu bergantung linear. Salah satu kebergantungan diperoleh
       dengan menetapkan $c_3=2$, yang memberikan $c_1=-2$ dan $c_2=-4$.
      \partsitem Tanpa perlu menyusun sistem, kita dapat melihat bahwa elemen
        kedua himpunan merupakan kelipatan elemen pertama (yakni, $0$ kali
        elemen pertama).
      \end{exparts}
    \end{answer}
  \recommended \item
    Buktikan bahwa masing-masing himpunan \( \set{f,g} \) berikut bebas linear
    dalam ruang vektor semua fungsi dari \( \Re^+ \) ke \( \Re \).
    \begin{exparts}
      \partsitem \( f(x)=x \) dan \( g(x)=1/x \)
      \partsitem \( f(x)=\cos(x) \) dan \( g(x)=\sin(x) \)
      \partsitem \( f(x)=e^x \) dan \( g(x)=\ln(x) \)
    \end{exparts}
    \begin{answer}
      Misalkan $\map{Z}{\Re^+}{\Re}$ adalah fungsi nol $Z(x)=0$, yang
      merupakan identitas aditif dalam ruang vektor yang sedang dibahas.
      \begin{exparts}
        \partsitem Himpunan ini bebas linear.
          Perhatikan \( c_1\cdot f(x)+c_2\cdot g(x)=Z(x) \). Menyubstitusikan
          \( x=1 \) dan \( x=2 \) menghasilkan sistem linear
          \begin{equation*}
            \begin{linsys}{2}
              c_1\cdot 1  &+  &c_2\cdot 1     &=  &0  \\
              c_1\cdot 2  &+  &c_2\cdot (1/2) &=  &0
            \end{linsys}
          \end{equation*}
          dengan solusi tunggal \( c_1=0 \), \( c_2=0 \).
        \partsitem Himpunan ini bebas linear.
          Perhatikan \( c_1\cdot f(x)+c_2\cdot g(x)=Z(x) \), lalu
          substitusikan \( x=\pi \) dan \( x=\pi/2 \) untuk memperoleh
          \begin{equation*}
            \begin{linsys}{2}
              c_1\cdot (-1)  &+  &c_2\cdot 0     &=  &0  \\
              c_1\cdot 0     &+  &c_2\cdot 1     &=  &0
            \end{linsys}
          \end{equation*}
          yang jelas memberikan \( c_1=0 \), \( c_2=0 \).
        \partsitem Himpunan ini juga bebas linear.
          Dengan memperhatikan \( c_1\cdot f(x)+c_2\cdot g(x)=Z(x) \) dan
          menyubstitusikan \( x=1 \) serta \( x=e \),
          \begin{equation*}
            \begin{linsys}{2}
              c_1\cdot e    &+  &c_2\cdot 0     &=  &0  \\
              c_1\cdot e^e  &+  &c_2\cdot 1     &=  &0
            \end{linsys}
          \end{equation*}
          diperoleh \( c_1=0 \) dan \( c_2=0 \).
      \end{exparts}  
     \end{answer}
  \recommended \item 
    Manakah dari subhimpunan ruang fungsi bernilai real dengan satu variabel
    real berikut yang bergantung linear dan manakah yang bebas linear?
    (Beberapa fungsi konstan telah disingkat; misalnya, pada butir pertama,
    `$2$' menyatakan fungsi konstan $f(x)=2$.)
    \begin{exparts*}
      \partsitem \( \set{2,4\sin^2(x),\cos^2(x)} \)
      \partsitem \( \set{1,\sin(x),\sin(2x)} \)
      \partsitem \( \set{x,\cos(x)} \)
      \partsitem \( \set{(1+x)^2,x^2+2x,3} \)
      \partsitem \( \set{\cos(2x),\sin^2(x),\cos^2(x)} \)
      \partsitem \( \set{0,x,x^2} \)
    \end{exparts*}
    \begin{answer}
      Dalam setiap kasus, jika himpunannya bebas, Anda harus membuktikannya;
      jika bergantung, Anda harus memperlihatkan suatu kebergantungan.
      \begin{exparts}
        \partsitem Himpunan ini bergantung.
          Relasi yang sudah dikenal $\sin^2(x)+\cos^2(x)=1$ menunjukkan bahwa
          $2=c_1\cdot(4\sin^2(x))+c_2\cdot(\cos^2(x))$ dipenuhi oleh
          $c_1=1/2$ dan $c_2=2$.
        \partsitem Himpunan ini bebas.
          Perhatikan relasi
          $c_1\cdot 1+c_2\cdot\sin(x)+c_3\cdot\sin(2x)=0$
          (di sini `$0$' adalah fungsi nol). Mengambil tiga titik yang sesuai,
          seperti $x=\pi$, $x=\pi/2$, $x=\pi/4$, menghasilkan sistem
          \begin{equation*}
            \begin{linsys}{3}
               c_1  &   &                &   &      &=  &0  \\
               c_1  &+  &c_2             &   &      &=  &0  \\
               c_1  &+  &(\sqrt{2}/2)c_2 &+  &c_3   &=  &0  
            \end{linsys}
          \end{equation*}
          yang satu-satunya solusinya ialah $c_1=0$, $c_2=0$, dan $c_3=0$.
        \partsitem Melalui pengamatan, himpunan ini bebas. Kebergantungan
          $\cos(x)=c\cdot x$ tidak mungkin ada karena fungsi kosinus bukan
          kelipatan fungsi identitas (kita menerapkan
          \nearbycorollary{cor:LDMeansLC}).
        \partsitem Melalui pengamatan, kita menemukan adanya kebergantungan.
          Karena $(1+x)^2=x^2+2x+1$, persamaan
          $c_1\cdot(1+x)^2+c_2\cdot(x^2+2x)=3$ dipenuhi oleh $c_1=3$ dan
          $c_2=-3$.
        \partsitem Himpunan ini bergantung. Cara termudah untuk melihatnya
          ialah mengingat relasi trigonometri
          $\cos^2(x)-\sin^2(x)=\cos(2x)$.
          (\textit{Catatan.} Orang yang tidak mengingat relasi ini dan mencoba
          beberapa nilai $x$ tidak akan pernah memperoleh sistem yang menuju
          ke solusi tunggal, sehingga tidak akan pernah dapat menyimpulkan
          bahwa himpunan tersebut bebas. Tentu saja, orang itu mungkin bertanya
          apakah ia hanya belum mencoba himpunan nilai $x$ yang tepat, tetapi
          beberapa percobaan akan membuat kebanyakan orang mencari suatu
          kebergantungan sebagai gantinya.)
        \partsitem Himpunan ini bergantung karena memuat objek nol dalam ruang
          vektor, yaitu polinomial nol.
      \end{exparts}  
     \end{answer}
  \item 
    Apakah persamaan \( \sin^2(x)/\cos^2(x)=\tan^2(x) \) menunjukkan bahwa
    himpunan fungsi \( \set{\sin^2(x),\cos^2(x),\tan^2(x)} \) merupakan
    subhimpunan bergantung linear dari himpunan semua fungsi bernilai real
    dengan domain interval \( (-\pi/2..\pi/2) \), yakni bilangan real antara
    \( -\pi/2 \) dan \( \pi/2 \)?
    \begin{answer}
      Tidak, persamaan itu bukan relasi linear. Bahkan, himpunan ini bebas,
      sebagaimana ditunjukkan oleh sistem yang timbul ketika \( x \) diambil
      sebesar \( 0 \), \( \pi/6 \), dan \( \pi/4 \).
    \end{answer}
  \item  
    Apakah bidang-$xy$, sebagai subhimpunan ruang vektor $\Re^3$, bebas linear?
    \begin{answer}
      Tidak. Berikut dua anggota bidang tersebut, dengan anggota kedua
      merupakan kelipatan anggota pertama.
      \begin{equation*}
        \colvec{1  \\ 0 \\ 0},\,\colvec{2 \\ 0 \\ 0}
      \end{equation*}
      (Alasan lain bahwa jawabannya ``tidak'' ialah vektor nol merupakan
      anggota bidang tersebut dan tidak ada himpunan yang memuat vektor nol
      yang bebas linear.)
    \end{answer}
  \recommended \item
    Tunjukkan bahwa baris-baris tak nol dari suatu matriks berbentuk eselon
    membentuk himpunan bebas linear.
    \begin{answer}
      Kita telah menunjukkan hal ini: Lema Kombinasi Linear dan akibatnya
      menyatakan bahwa, dalam matriks berbentuk eselon, tidak ada baris tak nol
      yang merupakan kombinasi linear dari baris-baris lainnya.
    \end{answer}
  \item
     \begin{exparts}
       \partsitem Tunjukkan bahwa jika himpunan
          \( \set{\vec{u},\vec{v},\vec{w}} \) bebas linear, maka himpunan
          \( \set{\vec{u},\vec{u}+\vec{v},\vec{u}+\vec{v}+\vec{w}} \)
          juga bebas linear.
       \partsitem Apa hubungan antara kebebasan atau kebergantungan linear
         \( \set{\vec{u},\vec{v},\vec{w}} \) dan kebebasan atau kebergantungan
         \( \set{\vec{u}-\vec{v},\vec{v}-\vec{w},\vec{w}-\vec{u}} \)?
     \end{exparts}
     \begin{answer}
       \begin{exparts}
         \partsitem Andaikan \( \set{\vec{u},\vec{v},\vec{w}} \) bebas
           linear, sehingga setiap relasi
           $d_0\vec{u}+d_1\vec{v}+d_2\vec{w}=\zero$ menghasilkan kesimpulan
           bahwa $d_0=0$, $d_1=0$, dan $d_2=0$.

           Perhatikan relasi
           \( c_1(\vec{u})+c_2(\vec{u}+\vec{v})+c_3(\vec{u}+\vec{v}+\vec{w})
           =\zero \).
           Tulis ulang relasi itu untuk memperoleh
           \( (c_1+c_2+c_3)\vec{u}+(c_2+c_3)\vec{v}+(c_3)\vec{w}=\zero \).
           Dengan mengambil $d_0$ sebesar $c_1+c_2+c_3$, $d_1$ sebesar
           $c_2+c_3$, dan $d_2$ sebesar $c_3$, kita memperoleh sistem berikut.
           \begin{equation*}
             \begin{linsys}{3}
               c_1  &+  &c_2  &+  &c_3  &=  &0  \\
                    &   &c_2  &+  &c_3  &=  &0  \\
                    &   &     &   &c_3  &=  &0
             \end{linsys}
           \end{equation*}
           Kesimpulannya: semua $c$ bernilai nol, sehingga himpunan tersebut
           bebas linear.
        \partsitem Himpunan kedua bergantung
           \begin{equation*}
             1\cdot(\vec{u}-\vec{v})
             +1\cdot(\vec{v}-\vec{w})
             +1\cdot(\vec{w}-\vec{u})
             =\zero
           \end{equation*}
           tanpa memandang apakah himpunan pertama bebas atau tidak.
      \end{exparts}
    \end{answer}
  \item 
    \nearbyexample{ex:EmSetLI} menunjukkan bahwa himpunan kosong bebas linear.
    \begin{exparts}
      \partsitem Kapan suatu himpunan beranggota satu bebas linear?
      \partsitem Bagaimana dengan himpunan beranggota dua?
    \end{exparts}
    \begin{answer}
      \begin{exparts}
        \partsitem Himpunan tunggal $\set{\vec{v}}$ bebas linear jika dan
          hanya jika $\vec{v}\neq\zero$. Untuk arah `jika', dengan
          $\vec{v}\neq\zero$, kita dapat menerapkan
          \nearbylemma{le:LDIffANonTrivLinRel} dengan memperhatikan relasi
          \( c\cdot\vec{v}=\zero \) dan mencatat bahwa satu-satunya solusi
          adalah solusi trivial:~$c=0$. Untuk arah `hanya jika', ingat saja
          bahwa
          \nearbyexample{ex:SetWithZeroVecLD}
          menunjukkan bahwa $\set{\zero}$ bergantung linear. Jadi, jika
          $\set{\vec{v}}$ bebas linear, maka $\vec{v}\neq\zero$.

          (\textit{Catatan.} Jawaban lain ialah mengatakan bahwa ini merupakan
          kasus khusus dari
          \nearbylemma{lm:AddVecLiSetIsLiIffVecNotInSpan} 
          ketika \( S=\emptyset \).)
        \partsitem Himpunan beranggota dua bebas linear jika dan hanya jika
          tidak satu pun anggotanya merupakan kelipatan anggota lainnya
          (perhatikan bahwa jika salah satunya adalah vektor nol, vektor itu
          merupakan kelipatan dari yang lain). Pernyataan ekuivalennya ialah:
          suatu himpunan bergantung linear jika dan hanya jika satu elemennya
          merupakan kelipatan elemen lainnya.

          Pembuktiannya mudah. Himpunan $\set{\vec{v}_1,\vec{v}_2}$ bergantung
          linear jika dan hanya jika terdapat relasi
          $c_1\vec{v}_1+c_2\vec{v}_2=\zero$ dengan $c_1\neq 0$ atau
          $c_2\neq 0$ (atau keduanya). Hal itu berlaku jika dan hanya jika
          $\vec{v}_1=(-c_2/c_1)\vec{v}_2$ atau
          $\vec{v}_2=(-c_1/c_2)\vec{v}_1$ (atau keduanya).
       \end{exparts}   
     \end{answer}
  \item  
    Dalam sebarang ruang vektor \( V \), himpunan kosong bebas linear.
    Bagaimana dengan seluruh \( V \)?
    \begin{answer}
      Himpunan ini bergantung linear karena memuat vektor nol.
    \end{answer}
  \item 
    Tunjukkan bahwa jika \( \set{\vec{x},\vec{y},\vec{z}} \) bebas linear,
    maka semua subhimpunan sejatinya juga bebas:~\( \set{\vec{x},\vec{y}} \),
    \( \set{\vec{x},\vec{z}} \), \( \set{\vec{y},\vec{z}} \),
    \( \set{\vec{x}} \),\( \set{\vec{y}} \), \( \set{\vec{z}} \),
    dan \( \set{} \).
    Apakah kebalikannya juga berlaku?
    \begin{answer}
      \nearbylemma{le:SubsetPreserveLI} memberikan arah `jika'. Kebalikannya
      (pernyataan `hanya jika') tidak berlaku. Sebagai contoh, perhatikan ruang
      vektor \( \Re^2 \) dan vektor-vektor berikut.
      \begin{equation*}
         \vec{x}=\colvec{1 \\ 0},\quad
         \vec{y}=\colvec{0 \\ 1},\quad
         \vec{z}=\colvec{1 \\ 1}
      \end{equation*} 
    \end{answer}
  \item 
    \begin{exparts}
      \partsitem Tunjukkan bahwa
        \begin{equation*}
          S=\set{\colvec{1 \\ 1 \\ 0},\colvec{-1 \\ 2 \\ 0}}
        \end{equation*}
        merupakan subhimpunan bebas linear dari \( \Re^3 \).
      \partsitem Tunjukkan bahwa
        \begin{equation*}
          \colvec{3 \\ 2 \\ 0}
        \end{equation*}
        berada dalam rentang $S$ dengan mencari \( c_1 \) dan \( c_2 \) yang
        menghasilkan relasi linear.
        \begin{equation*}
          c_1\colvec{1 \\ 1 \\ 0}
          +c_2\colvec{-1 \\ 2 \\ 0}
          =\colvec{3 \\ 2 \\ 0}
        \end{equation*}
        Tunjukkan bahwa pasangan \( c_1,c_2 \) tersebut tunggal.
      \partsitem Andaikan \( S \) merupakan subhimpunan suatu ruang vektor dan
        \( \vec{v} \) berada dalam \( \spanof{S} \), sehingga \( \vec{v} \)
        merupakan kombinasi linear vektor-vektor dalam \( S \). Buktikan bahwa
        jika \( S \) bebas linear, kombinasi linear vektor-vektor dalam
        \( S \) yang berjumlah \( \vec{v} \) bersifat tunggal (yakni, tunggal
        hingga pengurutan ulang dan penambahan atau pembuangan suku berbentuk
        \( 0\cdot\vec{s} \)). Jadi, \( S \) merupakan himpunan perentang
        minimal dalam arti kuat ini: setiap vektor dalam \( \spanof{S} \) merupakan
        kombinasi elemen-elemen $S$ sesedikit mungkin\Dash hanya satu kali.
      \partsitem
        Buktikan bahwa ketika \( S \) tidak bebas linear, dapat terjadi dua
        kombinasi linear berbeda berjumlah vektor yang sama.
    \end{exparts}
    \begin{answer}
      \begin{exparts}
        \partsitem Sistem linear yang muncul dari
          \begin{equation*}
            c_1\colvec{1 \\ 1 \\ 0}
            +c_2\colvec{-1 \\ 2 \\ 0}
            =\colvec{0 \\ 0 \\ 0}
          \end{equation*}
          memiliki solusi tunggal \( c_1=0 \) dan \( c_2=0 \).
        \partsitem Sistem linear yang muncul dari
          \begin{equation*}
            c_1\colvec{1 \\ 1 \\ 0}
            +c_2\colvec{-1 \\ 2 \\ 0}
            =\colvec{3 \\ 2 \\ 0}
          \end{equation*}
          memiliki solusi tunggal \( c_1=8/3 \) dan \( c_2=-1/3 \).
        \partsitem Andaikan \( S \) bebas linear. Andaikan kita memiliki
          $\vec{v}=c_1\vec{s}_1+\dots+c_n\vec{s}_n$ sekaligus
          $\vec{v}=d_1\vec{t}_1+\dots+d_m\vec{t}_m$ (dengan vektor-vektornya
          merupakan anggota $S$). Sekarang,
          \begin{equation*}
            c_1\vec{s}_1+\dots+c_n\vec{s}_n
            =\vec{v}
            =d_1\vec{t}_1+\dots+d_m\vec{t}_m
          \end{equation*}
          dapat ditulis ulang sebagai berikut.
          \begin{equation*}
            c_1\vec{s}_1+\dots+c_n\vec{s}_n
            -d_1\vec{t}_1-\dots-d_m\vec{t}_m
            =\zero
          \end{equation*}
          Mungkin beberapa $\vec{s}$ sama dengan beberapa $\vec{t}$; kita
          dapat menggabungkan koefisien-koefisien yang terkait (yakni, jika
          $\vec{s}_i=\vec{t}_j$, maka
          $\cdots+c_i\vec{s}_i+\dots-d_j\vec{t}_j-\cdots$ dapat ditulis ulang
          sebagai $\cdots+(c_i-d_j)\vec{s}_i+\cdots$). Persamaan itu merupakan
          relasi linear di antara anggota-anggota berbeda (setelah penggabungan
          selesai) dari himpunan $S$. Kita telah mengandaikan bahwa $S$ bebas
          linear, sehingga semua koefisiennya nol. Jika $i$ sedemikian sehingga
          $\vec{s}_i$ tidak sama dengan satu pun $\vec{t}_j$, maka $c_i$ nol.
          Jika $j$ sedemikian sehingga $\vec{t}_j$ tidak sama dengan satu pun
          $\vec{s}_i$, maka $d_j$ nol. Dalam kasus terakhir, kita memiliki
          $c_i-d_j=0$, sehingga $c_i=d_j$.

          Oleh karena itu, kedua jumlah semula sama, kecuali mungkin beberapa
          suku $0\cdot\vec{s}_i$ atau $0\cdot\vec{t}_j$ yang dapat kita abaikan.
        \partsitem
          Himpunan ini tidak bebas linear:
          \begin{equation*}
            S=\set{\colvec{1 \\ 0},\colvec{2 \\ 0}}\subset\Re^2
          \end{equation*}
          dan kedua kombinasi linear ini memberikan hasil yang sama
          \begin{equation*}
            \colvec{0 \\ 0}=2\cdot\colvec{1 \\ 0}-1\cdot\colvec{2 \\ 0}
                            =4\cdot\colvec{1 \\ 0}-2\cdot\colvec{2 \\ 0}
          \end{equation*}
          Jadi, suatu himpunan yang bergantung linear dapat memiliki dua
          kombinasi linear berbeda dengan hasil yang sama.

          Bahkan, pernyataan yang lebih kuat berikut berlaku: jika suatu
          himpunan bergantung linear, pasti terdapat dua kombinasi linear
          berbeda yang berjumlah vektor yang sama. Secara singkat, jika
          \( c_1\vec{s}_1+\dots+c_n\vec{s}_n=\zero \), mengalikan kedua ruas
          relasi tersebut dengan dua menghasilkan relasi lain. Jika relasi
          pertama tak trivial, relasi kedua juga tak trivial.
      \end{exparts}  
    \end{answer}
  \item  \label{exer:PolyZeroFcnOnlyIfZeroPol}
     Buktikan bahwa suatu polinomial menghasilkan fungsi nol jika dan hanya
     jika polinomial itu adalah polinomial nol.
     (\textit{Komentar.} Pertanyaan ini bukan masalah Aljabar Linear, tetapi
     hasilnya sering kita gunakan. Suatu polinomial menghasilkan fungsi secara
     alami:~$x\mapsto c_nx^n+\dots+c_1x+c_0$.)
     \begin{answer}
       Dalam pernyataan `jika dan hanya jika' ini, arah `jika' sudah
       jelas\Dash jika polinomialnya adalah polinomial nol, fungsi yang timbul
       dari polinomial itu pasti fungsi nol $x\mapsto 0$. Untuk arah `hanya
       jika', tulis $p(x)=c_nx^n+\dots+c_0$. Menyubstitusikan nol,
       $p(0)=0$, memberikan $c_0=0$. Mengambil turunannya dan menyubstitusikan
       nol, $p^\prime(0)=0$, memberikan $c_1=0$. Dengan cara serupa, kita
       memperoleh bahwa setiap $c_i$ nol dan $p$ adalah polinomial nol.
     \end{answer}
  \item
    Kembali ke Bagian 1.2 dan definisikan ulang titik, garis, bidang, serta
    permukaan linear lainnya agar kasus-kasus degenerat dihindari.
    \begin{answer}
      Pembahasan dalam bagian ini menyarankan agar kita mendefinisikan
      permukaan linear tak degenerat berdimensi \( n \) sebagai rentang suatu
      himpunan bebas linear yang terdiri atas \( n \) vektor.
    \end{answer}
  \item 
    \begin{exparts}  
      \partsitem Tunjukkan bahwa setiap himpunan empat vektor dalam \( \Re^2 \)
         bergantung linear.
      \partsitem Apakah ini berlaku untuk setiap himpunan lima vektor?
         Untuk setiap himpunan tiga vektor?
      \partsitem Berapa banyak elemen maksimum yang dapat dimiliki suatu
         subhimpunan bebas linear dari $\Re^2$?
    \end{exparts}
    \begin{answer}
      \begin{exparts}
        \partsitem Untuk sebarang $a_{1,1}$, \ldots, $a_{2,4}$,
          \begin{equation*}
            c_1\colvec{a_{1,1} \\ a_{2,1}}
            +c_2\colvec{a_{1,2} \\ a_{2,2}}
            +c_3\colvec{a_{1,3} \\ a_{2,3}}
            +c_4\colvec{a_{1,4} \\ a_{2,4}}
            =\colvec{0 \\ 0}
          \end{equation*}
          menghasilkan suatu sistem linear
          \begin{equation*}
             \begin{linsys}{4}
              a_{1,1}c_1 &+ &a_{1,2}c_2 &+ &a_{1,3}c_3 &+ &a_{1,4}c_4 &= &0  \\
              a_{2,1}c_1 &+ &a_{2,2}c_2 &+ &a_{2,3}c_3 &+ &a_{2,4}c_4 &= &0  
             \end{linsys}
          \end{equation*}
        menghasilkan sistem linear yang memiliki tak hingga banyak solusi
        (Metode Gauss menyisakan setidaknya dua variabel bebas). Oleh karena
        itu, terdapat relasi linear tak trivial di antara anggota-anggota
        $\Re^2$ yang diberikan.
      \partsitem Setiap himpunan lima vektor merupakan himpunan super dari
        suatu himpunan empat vektor, sehingga bergantung linear.

        Untuk tiga vektor dari $\Re^2$, argumen dari butir sebelumnya tetap
        berlaku, dengan sedikit perubahan bahwa Metode Gauss sekarang hanya
        menyisakan setidaknya satu variabel bebas (tetapi ini tetap memberikan
        tak hingga banyak solusi).
      \partsitem Butir sebelumnya menunjukkan bahwa tidak ada subhimpunan
        beranggota tiga dari $\Re^2$ yang bebas. Kita tahu bahwa ada
        subhimpunan beranggota dua dari $\Re^2$ yang bebas\Dash salah satunya
        adalah
        \begin{equation*}
          \set{\colvec{1  \\ 0},\colvec{0  \\ 1}}
        \end{equation*} 
        sehingga jawabannya adalah dua.
     \end{exparts}  
    \end{answer}
  \item
    Adakah himpunan empat vektor dalam \( \Re^3 \) sedemikian sehingga setiap
    tiga vektornya membentuk himpunan bebas linear?
    \begin{answer}
      Ada; berikut salah satunya.
      \begin{equation*}
         \set{\colvec{1 \\ 0 \\ 0},
              \colvec{0 \\ 1 \\ 0},
              \colvec{0 \\ 0 \\ 1},
              \colvec{1 \\ 1 \\ 1} }
      \end{equation*}  
    \end{answer}
  \item  
    Apakah setiap himpunan yang bergantung linear harus memiliki subhimpunan
    yang bergantung dan subhimpunan yang bebas?
    \begin{answer}
      Ya. Tetapkan sebarang ruang vektor~$V$, dan misalkan $S$ merupakan
      subhimpunan bergantung linear dari~$V$. Sebagai subhimpunan $S$ yang
      bergantung, kita dapat mengambil~$S$ sendiri. Sebagai subhimpunan $S$
      yang bebas, kita dapat mengambil himpunan kosong.
    \end{answer}
  \item  
    Dalam \( \Re^4 \), manakah himpunan bebas linear terbesar yang dapat Anda
    temukan? Yang terkecil? Manakah himpunan bergantung linear terbesar? Yang
    terkecil? (`Terbesar' dan `terkecil' berarti tidak ada himpunan super atau
    subhimpunan dengan sifat yang sama.)
    \begin{answer} 
      Dalam \( \Re^4 \), himpunan bebas linear terbesar memiliki empat
      vektor. Ada banyak contoh himpunan semacam itu; berikut salah satunya.
      \begin{equation*}
        \set{\colvec{1 \\ 0 \\ 0 \\ 0},
             \colvec{0 \\ 1 \\ 0 \\ 0},
             \colvec{0 \\ 0 \\ 1 \\ 0},
             \colvec{0 \\ 0 \\ 0 \\ 1}  }
      \end{equation*}
      Untuk melihat bahwa tidak ada himpunan lima vektor atau lebih yang dapat
      bebas, susun
      \begin{equation*}
             c_1\colvec{a_{1,1} \\ a_{2,1} \\ a_{3,1} \\ a_{4,1}}
            +c_2\colvec{a_{1,2} \\ a_{2,2} \\ a_{3,2} \\ a_{4,2}}
            +c_3\colvec{a_{1,3} \\ a_{2,3} \\ a_{3,3} \\ a_{4,3}}
            +c_4\colvec{a_{1,4} \\ a_{2,4} \\ a_{3,4} \\ a_{4,4}}
            +c_5\colvec{a_{1,5} \\ a_{2,5} \\ a_{3,5} \\ a_{4,5}}
             =\colvec{0 \\ 0 \\ 0 \\ 0}
      \end{equation*}
      dan perhatikan bahwa sistem linear yang dihasilkan
      \begin{equation*}
        \begin{linsys}{5}
          a_{1,1}c_1 &+ &a_{1,2}c_2 &+ &a_{1,3}c_3 
              &+ &a_{1,4}c_4 &+ &a_{1,5}c_5 &=  &0   \\
          a_{2,1}c_1 &+ &a_{2,2}c_2 &+ &a_{2,3}c_3 
              &+ &a_{2,4}c_4 &+ &a_{2,5}c_5 &=  &0   \\
          a_{3,1}c_1 &+ &a_{3,2}c_2 &+ &a_{3,3}c_3 
              &+ &a_{3,4}c_4 &+ &a_{3,5}c_5 &=  &0   \\
          a_{4,1}c_1 &+ &a_{4,2}c_2 &+ &a_{4,3}c_3 
              &+ &a_{4,4}c_4 &+ &a_{4,5}c_5 &=  &0     
        \end{linsys}
      \end{equation*}
      memiliki empat persamaan dan lima variabel tak diketahui. Karena itu,
      Metode Gauss harus berakhir dengan setidaknya satu variabel \( c \) bebas,
      sehingga terdapat tak hingga banyak solusi. Jadi, relasi linear di atas
      di antara vektor-vektor berkomponen empat memiliki solusi selain solusi
      trivial.

      Himpunan bebas linear terkecil adalah himpunan kosong.

      Himpunan bergantung linear terbesar adalah \( \Re^4 \). Yang terkecil
      adalah \( \set{\zero} \).
    \end{answer}
  \recommended \item 
    Kebebasan linear dan kebergantungan linear merupakan sifat himpunan.
    Karena itu, secara alami kita dapat bertanya bagaimana sifat kebebasan dan
    kebergantungan linear berperilaku terhadap relasi dan operasi himpunan
    elementer yang sudah dikenal. Dalam isi subbagian ini, kita telah membahas
    relasi subhimpunan dan himpunan super. Kita juga dapat mempertimbangkan
    operasi irisan, komplemen, dan gabungan.
    \begin{exparts}
       \partsitem Bagaimana kebebasan linear berkaitan dengan irisan: dapatkah
         irisan dua himpunan bebas linear bersifat bebas? Haruskah demikian?
       \partsitem Bagaimana kebebasan linear berkaitan dengan komplemen?
       \partsitem Tunjukkan bahwa gabungan dua himpunan bebas linear dapat
          bersifat bebas linear.
       \partsitem Tunjukkan bahwa gabungan dua himpunan bebas linear tidak
          selalu bebas linear.
    \end{exparts}
    \begin{answer}
      \begin{exparts}
        \partsitem Irisan dua himpunan bebas linear $S\intersection T$ harus
          bebas linear karena merupakan subhimpunan dari himpunan bebas linear
          $S$ (dan tentu saja juga dari himpunan bebas linear $T$).
        \partsitem Komplemen suatu himpunan bebas linear bergantung linear
          karena memuat vektor nol.
        \partsitem Contoh sederhana dalam \( \Re^2 \) ialah kedua himpunan ini.
          \begin{equation*}
            S=\set{\colvec{1 \\ 0}}
            \qquad
            T=\set{\colvec{0 \\ 1}}          
          \end{equation*}
          Contoh yang sedikit lebih halus, juga dalam \( \Re^2 \), ialah kedua
          himpunan ini.
          \begin{equation*}
            S=\set{\colvec{1 \\ 0}}
            \qquad
            T=\set{\colvec{1 \\ 0}, \colvec{0 \\ 1}}
          \end{equation*}
        \partsitem Kita harus memberikan suatu contoh. Salah satu contoh dalam
          \( \Re^2 \) adalah
          \begin{equation*}
            S=\set{\colvec{1 \\ 0}}
            \qquad
            T=\set{\colvec{2 \\ 0}}
          \end{equation*}
          karena kebergantungan linear \( S\union T \) mudah dilihat.
      \end{exparts}
    \end{answer}
  \item \textit{Lanjutan dari latihan sebelumnya.}
    Bagaimana interaksi antara sifat kebebasan linear dan operasi gabungan?
    \begin{exparts}
       \partsitem Kita mungkin menduga bahwa gabungan $S\union T$ dari dua
          himpunan bebas linear bersifat bebas linear jika dan hanya jika
          rentang keduanya memiliki irisan trivial
          $\spanof{S}\intersection \spanof{T}=\set{\zero}$. Apa yang salah
          dengan argumen berikut untuk arah `hanya jika' dari dugaan tersebut?
          ``Jika gabungan $S\union T$ bebas linear, satu-satunya solusi dari
          $c_1\vec{s}_1+\cdots+c_n\vec{s}_n
            +d_1\vec{t}_1+\cdots+d_m\vec{t}_m
            =\zero$
          adalah solusi trivial $c_1=0$, \ldots, $d_m=0$. Jadi, setiap anggota
          irisan rentang harus merupakan vektor nol karena dalam
          $c_1\vec{s}_1+\cdots+c_n\vec{s}_n
            =d_1\vec{t}_1+\cdots+d_m\vec{t}_m$
          setiap skalar bernilai nol.''
       \partsitem Berikan contoh yang menunjukkan bahwa dugaan tersebut salah.
       \partsitem Carilah himpunan bebas linear \( S \) dan \( T \) sedemikian
          sehingga gabungan \( S-(S\cap T)\) dan \( T-(S\cap T) \) bebas
          linear, tetapi gabungan \( S\cup T \) tidak bebas linear.
       \partsitem Karakterisasikan kapan gabungan dua himpunan bebas linear
          bersifat bebas linear dalam suku-suku irisan rentangnya.
    \end{exparts}
    \begin{answer}
      \begin{exparts}
        \partsitem \nearbylemma{le:LDIffANonTrivLinRel} mengharuskan
          vektor-vektor
          $\vec{s}_1,\ldots,\vec{s}_n,\vec{t}_1,\ldots,\vec{t}_m$ berbeda.
          Namun, gabungan $S\cup T$ dapat bebas linear dengan suatu
          $\vec{s}_i$ sama dengan suatu $\vec{t}_j$.
        \partsitem Salah satu contoh dalam \( \Re^2 \) ialah kedua himpunan ini.
          \begin{equation*}
            S=\set{\colvec{1 \\ 0}}
            \qquad
            T=\set{\colvec{1 \\ 0}, \colvec{0 \\ 1}}
          \end{equation*}
        \partsitem 
          Contoh dari \( \Re^2 \) ialah himpunan-himpunan berikut.
          \begin{equation*}
            S=\set{\colvec{1 \\ 0}, \colvec{0 \\ 1}}
            \qquad
            T=\set{\colvec{1 \\ 0}, \colvec{1 \\ 1}}
          \end{equation*}
        \partsitem Gabungan dua himpunan bebas linear $S\union T$ bebas linear
          jika dan hanya jika rentang \( S \) dan \( T-(S\cap T)\) memiliki
          irisan trivial
          $\spanof{S}\intersection \spanof{T-(S\cap T)}=\set{\zero}$.
          Untuk membuktikannya, andaikan \( S \) dan \( T \) merupakan
          subhimpunan bebas linear dari suatu ruang vektor.

          Untuk arah `jika', andaikan irisan rentang tersebut trivial,
          \( \spanof{S}\intersection \spanof{T-(S\cap T)}=\set{\zero} \).
          Perhatikan himpunan $S\union (T-(S\cap T))=S\union T$ dan relasi
          linear
          $c_1\vec{s}_1+\dots+c_n\vec{s}_n
            +d_1\vec{t}_1+\dots+d_m\vec{t}_m=\zero$.
          Pengurangan menghasilkan
          $c_1\vec{s}_1+\dots+c_n\vec{s}_n=
            -d_1\vec{t}_1-\dots-d_m\vec{t}_m$.
          Ruas kiri persamaan itu berjumlah suatu vektor dalam $\spanof{S}$,
          sedangkan ruas kanannya merupakan vektor dalam
          $\spanof{T-(S\cap T)}$. Oleh karena irisan rentangnya trivial, kedua
          ruas sama dengan vektor nol. Karena $S$ bebas linear, semua $c$
          bernilai nol. Karena $T$ bebas linear, $T-(S\cap T)$ juga bebas
          linear, sehingga semua $d$ bernilai nol. Jadi, relasi linear semula
          di antara anggota $S\union T$ hanya berlaku jika semua koefisiennya
          nol. Dengan demikian, $S\union T$ bebas linear.

          Untuk arah `hanya jika', kita dapat menjalankan argumen yang sama secara
          terbalik. Andaikan gabungan $S\union T$ bebas linear. Perhatikan
          suatu relasi linear di antara anggota $S$ dan $T-(S\cap T)$:
          $c_1\vec{s}_1+\cdots+c_n\vec{s}_n
            +d_1\vec{t}_1+\cdots+d_m\vec{t}_m
            =\zero$
          Perhatikan bahwa tidak ada $\vec{s}_i$ yang sama dengan
          $\vec{t}_j$, sehingga ini merupakan kombinasi vektor-vektor berbeda,
          sebagaimana disyaratkan \nearbylemma{le:LDIffANonTrivLinRel}. Jadi,
          satu-satunya solusi adalah solusi trivial $c_1=0$, \ldots, $d_m=0$.
          Karena setiap vektor $\vec{v}$ dalam irisan rentang
          $\spanof{S}\intersection \spanof{T-(S\cap T)}$ dapat kita tulis
          sebagai
          $\vec{v}=c_1\vec{s}_1+\cdots+c_n\vec{s}_n
            =-d_1\vec{t}_1-\cdots-d_m\vec{t}_m$,
          vektor itu harus merupakan vektor nol karena setiap skalarnya nol.
      \end{exparts}  
    \end{answer}
  \item \label{exer:FillIndDetProofSetHasLISub}
    Untuk \nearbycorollary{th:AlwaysAnLDSubset},
    \begin{exparts}
       \partsitem lengkapilah induksi dalam pembuktiannya;
       \partsitem berikan pembuktian alternatif yang dimulai dengan himpunan
         kosong dan membangun suatu barisan subhimpunan bebas linear dari
         himpunan hingga yang diberikan sampai muncul subhimpunan dengan
         rentang yang sama seperti himpunan tersebut.
    \end{exparts}
    \begin{answer}
      \begin{exparts}
         \partsitem Kita melakukan induksi pada banyaknya vektor dalam
           himpunan hingga \( S \).

           Kasus dasarnya ialah $S$ tidak memiliki elemen. Dalam kasus ini,
           $S$ bebas linear dan tidak ada yang perlu diperiksa\Dash subhimpunan
           $S$ yang memiliki rentang sama dengan $S$ adalah $S$ sendiri.

           Untuk langkah induktif, andaikan teorema itu benar untuk semua
           himpunan berukuran $n=0$, $n=1$, \ldots, $n=k$, guna membuktikan
           bahwa teorema berlaku ketika \( S \) memiliki $n=k+1$ elemen. Jika
           himpunan beranggota $k+1$
           \( S=\set{\vec{s}_0,\dots,\vec{s}_{k}} \) bebas linear, teorema ini
           trivial; jadi, andaikan himpunan tersebut bergantung. Berdasarkan
           \nearbycorollary{cor:LDMeansLC}, terdapat suatu \( \vec{s}_i \)
           yang merupakan kombinasi linear vektor-vektor lain dalam \( S \).
           Definisikan \( S_1=S-\set{\vec{s}_i} \), dan perhatikan bahwa
           \( S_1 \) memiliki rentang yang sama dengan \( S \) berdasarkan
           \nearbycorollary{cor:OmitVecNotShrinkSpanIffVecIsDep}. Himpunan
           \( S_1 \) memiliki \( k \) elemen, sehingga hipotesis induktif
           berlaku dan memberikan subhimpunan bebas linear dengan rentang yang
           sama. Subhimpunan dari \( S_1 \) itu adalah subhimpunan \( S \) yang
           diinginkan.
         \partsitem Berikut sketsa argumennya. Kita menghilangkan rincian
           argumen induksi.

           Jika himpunan hingga \( S \) kosong, tidak ada yang perlu
           dibuktikan. Jika \( S=\set{\zero} \), subhimpunan kosong memenuhi.

           Jika tidak, ambil suatu vektor tak nol \( \vec{s}_1\in S \) dan
           definisikan \( S_1=\set{\vec{s}_1} \). Jika
           \( \spanof{S_1}=\spanof{S} \), pembuktian selesai dengan mencatat
           bahwa \( S_1 \) bebas linear.

           Jika tidak, terdapat suatu vektor tak nol
           \( \vec{s}_2\in S-\spanof{S_1} \) (jika setiap \( \vec{s}\in S \)
           berada dalam \( \spanof{S_1} \), maka
           \( \spanof{S_1}=\spanof{S} \)). Definisikan
           \( S_2=S_1\union\set{\vec{s}_2} \). Jika
           \( \spanof{S_2}=\spanof{S} \), pembuktian selesai dengan memakai
           \nearbycorollary{cor:LDMeansLC} untuk menunjukkan bahwa \( S_2 \)
           bebas linear.

           Ulangi paragraf terakhir sampai muncul himpunan dengan rentang yang
           cukup besar. Hal itu pada akhirnya pasti terjadi karena \( S \)
           hingga, dan paling lambat \( \spanof{S} \) akan tercapai ketika kita
           telah memakai setiap vektor dari \( S \).
      \end{exparts}  
    \end{answer}
  \item 
     Dengan sejumlah perhitungan, kita dapat memperoleh rumus untuk menentukan
     apakah suatu himpunan vektor bebas linear atau tidak.
     \begin{exparts}
       \partsitem Tunjukkan bahwa subhimpunan \( \Re^2 \) berikut
         \begin{equation*}
           \set{\colvec{a \\ c},\colvec{b \\ d}}
         \end{equation*}
         bebas linear jika dan hanya jika \( ad-bc\neq 0 \).
       \partsitem Tunjukkan bahwa subhimpunan \( \Re^3 \) berikut
         \begin{equation*}
           \set{\colvec{a \\ d \\ g},
                \colvec{b \\ e \\ h},
                \colvec{c \\ f \\ i}  }
         \end{equation*}
         bebas linear jika dan hanya jika
         \( aei+bfg+cdh-hfa-idb-gec \neq 0 \).
       \partsitem Kapan subhimpunan \( \Re^3 \) berikut
         \begin{equation*}
           \set{\colvec{a \\ d \\ g},
                \colvec{b \\ e \\ h} }
         \end{equation*}
         bebas linear?
       \partsitem Ini adalah pertanyaan pendapat: untuk himpunan empat vektor
         dari \( \Re^4 \), haruskah ada suatu rumus yang melibatkan keenam
         belas entri dan menentukan kebebasan himpunan tersebut? (Anda tidak
         perlu menghasilkan rumus semacam itu; cukup putuskan apakah rumus itu
         ada.)
    \end{exparts}
    \begin{answer}
      \begin{exparts}
        \partsitem 
          Dengan terlebih dahulu mengandaikan \( a\neq 0 \),
          \begin{equation*}
            x\colvec{a \\ c}
            +y\colvec{b \\ d}
            =\colvec{0 \\ 0}
          \end{equation*}
          menghasilkan
          \begin{equation*}
            \begin{linsys}{2}
               ax  &+  &by &=  &0  \\
               cx  &+  &dy &=  &0  
            \end{linsys}
            \grstep{-(c/a)\rho_1+\rho_2}
            \begin{linsys}{2}
               ax  &+  &by           &=  &0  \\
                   &   &(-(c/a)b+d)y &=  &0  
             \end{linsys}
          \end{equation*}
          yang hanya memiliki solusi trivial jika dan hanya jika
          \( 0\neq-(c/a)b+d=(-cb+ad)/a \)
          (dalam kasus ini kita telah mengandaikan \( a\neq 0 \), sehingga
          substitusi balik menghasilkan solusi tunggal).

          Kasus \( a=0 \) juga tidak sulit\Dash pisahkan menjadi subkasus
          \( c\neq 0 \) dan \( c=0 \), lalu perhatikan bahwa dalam kasus-kasus
          ini \( ad-bc=0\cdot d-bc \).

          \textit{Komentar.} Latihan sebelumnya menunjukkan bahwa himpunan dua
          vektor bergantung linear jika dan hanya jika salah satu vektor
          merupakan kelipatan skalar dari vektor lainnya. Kita juga dapat
          memakainya untuk melakukan perhitungan.
        \partsitem Persamaan
          \begin{equation*}
            c_1\colvec{a \\ d \\ g}
            +c_2\colvec{b \\ e \\ h}
            +c_3\colvec{c \\ f \\ i}
            =\colvec{0 \\ 0 \\ 0}
          \end{equation*}
         menyatakan suatu sistem linear homogen. Kita melanjutkan dengan
         menuliskannya dalam bentuk matriks dan menerapkan Metode Gauss.

         Pertama-tama kita reduksi matriks menjadi segitiga atas. Andaikan
         \( a\neq 0 \). Dengan itu, kita dapat menolkan entri-entri di bawah
         entri pertama pada kolom pertama.
         \begin{equation*}
           \grstep{(1/a)\rho_1}
           \begin{amat}{3}
              1   &b/a   &c/a  &0 \\
              d   &e     &f    &0 \\
              g   &h     &i    &0
            \end{amat}                                            
           \grstep[-g\rho_1+\rho_3]{-d\rho_1+\rho_2}
           \begin{amat}{3}
              1   &b/a           &c/a        &0   \\
              0   &(ae-bd)/a     &(af-cd)/a  &0   \\
              0   &(ah-bg)/a     &(ai-cg)/a  &0
            \end{amat}                                            
         \end{equation*}
         Kemudian kita memperoleh $1$ pada entri baris kedua, kolom kedua.
         (Untuk sementara, andaikan \( ae-bd\neq 0 \) agar langkah reduksi
         baris ini dapat dilakukan.)
         \begin{equation*}
           \grstep{(a/(ae-bd))\rho_2}
           \begin{amat}{3}
              1   &b/a           &c/a             &0  \\
              0   &1             &(af-cd)/(ae-bd) &0  \\
              0   &(ah-bg)/a     &(ai-cg)/a       &0
            \end{amat}
         \end{equation*}
         Kemudian, berdasarkan pengandaian tersebut, kita melakukan operasi
         baris $-((ah-bg)/a)\rho_2+\rho_3$ untuk memperoleh berikut ini.
         \begin{equation*}
           % \grstep{((ah-bg)/a)\rho_2+\rho_3}
           \begin{amat}{3}
              1   &b/a   &c/a                              &0 \\
              0   &1     &(af-cd)/(ae-bd)                   &0 \\
              0   &0     &(aei+bgf+cdh-hfa-idb-gec)/(ae-bd) &0
            \end{amat}
         \end{equation*}
         Oleh karena itu, sistem semula tak singular jika dan hanya jika entri
         \( 3,3 \) di atas tak nol (pecahan ini terdefinisi karena pengandaian
         \( ae-bd\neq 0 \)). Entri tersebut sama dengan nol jika dan hanya jika
         pembilangnya nol.

         Berikutnya kita tangani pengandaian-pengandaian tersebut. Pertama,
         jika \( a\neq 0 \) tetapi \( ae-bd=0 \), kita mempertukarkan baris
         \begin{multline*}
           \begin{amat}{3}
              1   &b/a           &c/a        &0   \\
              0   &0             &(af-cd)/a  &0   \\
              0   &(ah-bg)/a     &(ai-cg)/a  &0
            \end{amat}                                    \\
           \grstep{\rho_2\leftrightarrow\rho_3}
           \begin{amat}{3}
              1   &b/a           &c/a        &0   \\
              0   &(ah-bg)/a     &(ai-cg)/a  &0   \\
              0   &0             &(af-cd)/a  &0
            \end{amat}
         \end{multline*}
         dan menyimpulkan bahwa sistem tersebut tak singular jika dan hanya
         jika \( ah-bg\neq 0 \) dan \( af-cd\neq 0 \). Ini sama dengan
         mensyaratkan agar hasil kali keduanya tak nol:
         \begin{align*}
            ahaf-ahcd-bgaf+bgcd
            &\neq 0                   \\
            ahaf-ahcd-bgaf+aegc
            &\neq 0                   \\
            a(haf-hcd-bgf+egc)
            &\neq 0
         \end{align*}
         (dalam peralihan dari baris pertama ke baris kedua, kita menerapkan
         pengandaian kasus $ae-bd=0$ dengan mengganti $bd$ oleh $ae$). Karena
         kita mengandaikan \( a\neq 0 \), kita memperoleh
         \( haf-hcd-bgf+egc\neq 0 \). Dengan $ae-bd=0$, kita dapat menulis
         ulang ini dalam bentuk yang diperlukan: dalam kasus \( a\neq 0 \) dan
         \( ae-bd=0 \) ini, sistem yang diberikan tak singular ketika
         \( haf-hcd-bgf+egc-i(ae-bd)\neq 0 \), sebagaimana diperlukan.

         Kasus-kasus sisanya memiliki sifat yang sama. Tangani kasus \( a=0 \)
         tetapi \( d\neq 0 \), serta kasus \( a=0 \), \( d=0 \), tetapi
         \( g\neq 0 \), dengan terlebih dahulu mempertukarkan baris lalu
         melanjutkan seperti di atas. Kasus \( a=0 \), \( d=0 \), dan
         \( g=0 \) mudah\Dash himpunan yang memuat vektor nol bergantung linear,
         dan rumusnya bernilai nol.
       \partsitem Himpunan itu bergantung linear jika dan hanya jika salah
         satu vektor merupakan kelipatan vektor lainnya. Artinya, himpunan itu
         tidak bebas jika dan hanya jika
         \begin{equation*}
           \colvec{a \\ d \\ g}=r\cdot\colvec{b \\ e \\ h}
          \quad\text{atau}\quad
           \colvec{b \\ e \\ h}=s\cdot\colvec{a \\ d \\ g}
         \end{equation*}
         (atau keduanya) untuk sejumlah skalar $r$ dan $s$. Dengan
         mengeliminasi $r$ dan $s$ agar syarat ini dinyatakan hanya dalam
         huruf-huruf yang diberikan, $a$, $b$, $d$, $e$, $g$, $h$, kita
         memperoleh bahwa himpunan itu tidak bebas\Dash yakni bergantung\Dash
         jika dan hanya jika \( ae-bd=ah-gb=dh-ge=0 \).
       \partsitem Kebergantungan atau kebebasan merupakan fungsi dari
         entri-entrinya, sehingga memang terdapat sebuah rumus (walaupun pada
         pandangan pertama orang mungkin mengira bahwa rumus tersebut
         melibatkan beberapa kasus: ``jika komponen pertama dari vektor
         pertama nol, maka \ldots'', dugaan ini ternyata tidak benar).
      \end{exparts}  
    \end{answer}
  \recommended \item  
    \begin{exparts}
      \partsitem Buktikan bahwa himpunan dua vektor tak nol yang saling tegak
        lurus dari \( \Re^n \) bebas linear ketika \( n>1 \).
      \partsitem Bagaimana jika \( n=1 \)? Jika \( n=0 \)?
      \partsitem Generalisasikan untuk lebih dari dua vektor.
    \end{exparts}
    \begin{answer}
      Ingat bahwa dua vektor dari \( \Re^n \) saling tegak lurus jika dan hanya
      jika hasil kali titik keduanya nol.
      \begin{exparts}
         \partsitem Andaikan \( \vec{v} \) dan \( \vec{w} \) merupakan vektor
           tak nol yang saling tegak lurus dalam $\Re^n$, dengan \( n>1 \).
           Pada relasi linear \( c\vec{v}+d\vec{w}=\zero \), ambil hasil kali
           titik kedua ruas dengan \( \vec{v} \) untuk menyimpulkan bahwa
           \( c\cdot\norm{\vec{v}}^2+d\cdot 0=0 \). Karena
           \( \vec{v}\neq\zero \), kita memperoleh \( c=0 \). Penerapan serupa
           dengan \( \vec{w} \) menunjukkan bahwa \( d=0 \).
         \partsitem Dua vektor dalam \( \Re^1 \) saling tegak lurus jika dan
           hanya jika setidaknya salah satunya nol.

           Kita mendefinisikan \( \Re^0 \) sebagai ruang trivial, sehingga
           $\vec{v}$ dan $\vec{w}$ keduanya merupakan vektor nol.
         \partsitem Generalisasi yang tepat ialah meninjau himpunan vektor
           \( \set{\vec{v}_1,\dots,\vec{v}_n}\subseteq\Re^k \) yang
           \definend{saling ortogonal} (juga disebut
           \definend{saling tegak lurus berpasangan}): jika \( i\neq j \), maka
           \( \vec{v}_i \) tegak lurus terhadap \( \vec{v}_j \). Meniru
           pembuktian butir pertama di atas menunjukkan bahwa himpunan
           vektor-vektor tak nol semacam itu bebas linear.
      \end{exparts}  
    \end{answer}
  \item 
    Perhatikan himpunan fungsi dari interval
    $\openinterval{-1}{1}\subseteq \Re$ ke $\Re$.
    \begin{exparts}
      \partsitem Tunjukkan bahwa himpunan ini merupakan ruang vektor di bawah
         operasi-operasi biasa.
      \partsitem Ingat rumus jumlah deret geometri tak hingga:
         \( 1+x+x^2+\cdots=1/(1-x) \) untuk semua \( x\in(-1..1) \). Mengapa
         ini tidak menyatakan suatu kebergantungan di dalam himpunan
         $\set{g(x)=1/(1-x),f_0(x)=1,f_1(x)=x,f_2(x)=x^2,\ldots}$ (dalam ruang
         vektor yang sedang kita pertimbangkan)?
         (\textit{Petunjuk.} Tinjau kembali definisi kombinasi linear.)
      \partsitem Tunjukkan bahwa himpunan dalam butir sebelumnya bebas linear.
    \end{exparts}
    Ini menunjukkan bahwa ada ruang vektor yang memiliki subhimpunan bebas
    linear tak hingga.
    \begin{answer}
      \begin{exparts}
        \partsitem Pemeriksaan ini rutin.
        \partsitem Penjumlahannya tak hingga (memiliki tak hingga banyak suku).
          Definisi kombinasi linear hanya melibatkan jumlah hingga.
        \partsitem Tidak ada jumlah hingga tak trivial dari anggota
           \( \set{g,f_0,f_1,\ldots} \) yang berjumlah objek nol: andaikan
           \begin{equation*}
              c_0\cdot (1/(1-x))+c_1\cdot 1+\dots+c_n\cdot x^n=0
           \end{equation*}
           (setiap jumlah hingga memakai pangkat tertinggi, di sini \( n \)).
           Kalikan kedua ruas dengan \( 1-x \) untuk menyimpulkan bahwa setiap
           koefisien nol, karena suatu polinomial mendeskripsikan fungsi nol
           hanya ketika polinomial itu adalah polinomial nol.
      \end{exparts}
     \end{answer}
  \item 
    Tunjukkan bahwa, jika \( S \) merupakan subruang dari \( V \), dan suatu
    subhimpunan $T$ dari \( S \) bebas linear dalam \( S \), maka $T$ juga
    bebas linear dalam \( V \). Apakah kebalikannya berlaku?
    \begin{answer}
      Kedua arah, `jika' dan `hanya jika', berlaku.

      Misalkan \( T \) merupakan subhimpunan dari subruang \( S \) pada ruang
      vektor \( V \). Pernyataan bahwa setiap relasi linear
      $c_1\vec{t}_1+\dots+c_n\vec{t}_n=\zero$ di antara anggota \( T \) harus
      merupakan relasi trivial $c_1=0$, \ldots, $c_n=0$ berlaku dalam
      \( S \) jika dan hanya jika berlaku dalam \( V \), karena subruang
      \( S \) mewarisi operasi penjumlahan dan perkalian skalarnya dari
      \( V \).
    \end{answer}
\end{exercises}
